\section{未锚定账本事件的叙事补充：eastern-han}

\label{sec:anchor-batch-two-eastern-han}

以下事件均为此前正文尚未设置导航目标的账本记录。每一锚点置于来源、制度与地方条件的叙事中，避免把年表、政权分类或战报修辞当作社会全貌。

\subsection{制度与地方资源}

\eventanchor{EH-0184-YELLOW-TURBAN}{中平元年二月（184年）·张角集团发动黄巾起事，冀、豫、兖等州郡响应，朝廷全国征调平叛}账本记录东汉与黄巾军在中平元年二月（184年）发生“张角集团发动黄巾起事，冀、豫、兖等州郡响应，朝廷全国征调平叛”。条目把背景概括为太平道网络在张角计划泄露后提前举事，地方积弊和灾荒使基层动员迅速扩散，并指出中央倚赖将领、豪强和州郡募兵应战，常备官僚秩序开始让位于地方军事资源。这个节点进入区域社会时，要经由文书、道路、军队、市场、寺院和家庭等不同中介；因此，政治行动的宣布、地方接收和日常生活的改变不能被视作同一时刻完成。

本条依据《后汉书·灵帝纪》；《后汉书·皇甫嵩传》；《后汉纪》卷一百零九，证据提示为“起事月份可靠；参与人数、同时性和各地组织联系有争议”。这些材料可支持年序、官方决定、战事或特定地点的线索，却难以直接统计所有人的财富、迁徙、信仰和动机。更稳妥的读法是区分诏令与实施、军报与人口、行政称谓与自我认同，并让地方材料的缺环限制解释的范围。

由此可见，该事件不是孤立的年表点，而是资源、信息与权力关系重新配置的一环。不同地区的官员、社群和家庭可能以不同方式承担征发、保护、迁移或宗教活动；现存来源只能照亮其中一部分，不能替沉默者制造统一的经验。

\eventanchor{EH-0184-WANCHENG-YELLOW-TURBAN}{中平元年六月（184年）·朱儁等攻破宛城黄巾军，南阳战区主力受挫}账本记录东汉与荆州黄巾军在中平元年六月（184年）发生“朱儁等攻破宛城黄巾军，南阳战区主力受挫”。条目把背景概括为黄巾在南阳依托人口与交通节点聚集，朝廷需切断其与其他州郡的联络，并指出荆州黄巾主力被击散，但残余武装和地方自保武力持续存在，平乱并未恢复原有行政秩序。这个节点进入区域社会时，要经由文书、道路、军队、市场、寺院和家庭等不同中介；因此，政治行动的宣布、地方接收和日常生活的改变不能被视作同一时刻完成。

本条依据《后汉书·朱儁传》；《后汉书·皇甫嵩传》；《资治通鉴》卷五十八，证据提示为“破城时间大体可靠；围城战术、斩获与黄巾残部去向有争议”。这些材料可支持年序、官方决定、战事或特定地点的线索，却难以直接统计所有人的财富、迁徙、信仰和动机。更稳妥的读法是区分诏令与实施、军报与人口、行政称谓与自我认同，并让地方材料的缺环限制解释的范围。

由此可见，该事件不是孤立的年表点，而是资源、信息与权力关系重新配置的一环。不同地区的官员、社群和家庭可能以不同方式承担征发、保护、迁移或宗教活动；现存来源只能照亮其中一部分，不能替沉默者制造统一的经验。

\eventanchor{EH-0184-GUANGZONG-QUYANG}{中平元年秋冬（184年）·皇甫嵩、卢植等在广宗、下曲阳击破张梁、张宝部，冀州黄巾主力覆灭}账本记录东汉与冀州黄巾军在中平元年秋冬（184年）发生“皇甫嵩、卢植等在广宗、下曲阳击破张梁、张宝部，冀州黄巾主力覆灭”。条目把背景概括为朝廷集中北方军队围攻黄巾据点，张角病死后起事军失去象征性核心，并指出黄巾的集中战场失败，皇甫嵩等军功上升；分散的黄巾和地方武装仍将影响后续军阀格局。这个节点进入区域社会时，要经由文书、道路、军队、市场、寺院和家庭等不同中介；因此，政治行动的宣布、地方接收和日常生活的改变不能被视作同一时刻完成。

本条依据《后汉书·皇甫嵩传》；《后汉书·卢植传》；《后汉纪》卷一百零九，证据提示为“战役结果可靠；卢植获罪、指挥分工和死伤数字有争议”。这些材料可支持年序、官方决定、战事或特定地点的线索，却难以直接统计所有人的财富、迁徙、信仰和动机。更稳妥的读法是区分诏令与实施、军报与人口、行政称谓与自我认同，并让地方材料的缺环限制解释的范围。

由此可见，该事件不是孤立的年表点，而是资源、信息与权力关系重新配置的一环。不同地区的官员、社群和家庭可能以不同方式承担征发、保护、迁移或宗教活动；现存来源只能照亮其中一部分，不能替沉默者制造统一的经验。

\eventanchor{EH-0184-LIANGZHOU-REBELLION}{中平元年（184年）·凉州军民与羌胡部众起事，北宫伯玉、边章、韩遂等先后成为叛军领袖}账本记录东汉、凉州叛军与羌胡集团在中平元年（184年）发生“凉州军民与羌胡部众起事，北宫伯玉、边章、韩遂等先后成为叛军领袖”。条目把背景概括为长期边防征发、羌乱遗绪和凉州军人利益冲突在黄巾危机中同时爆发，并指出西北形成长期自主军事集团，中央需委任董卓等带兵镇压，地方军权进一步制度化。这个节点进入区域社会时，要经由文书、道路、军队、市场、寺院和家庭等不同中介；因此，政治行动的宣布、地方接收和日常生活的改变不能被视作同一时刻完成。

本条依据《后汉书·灵帝纪》；《后汉书·董卓传》；《后汉书·西羌传》，证据提示为“起事年份可靠；领袖更替、族属称谓和各部关系有争议”。这些材料可支持年序、官方决定、战事或特定地点的线索，却难以直接统计所有人的财富、迁徙、信仰和动机。更稳妥的读法是区分诏令与实施、军报与人口、行政称谓与自我认同，并让地方材料的缺环限制解释的范围。

由此可见，该事件不是孤立的年表点，而是资源、信息与权力关系重新配置的一环。不同地区的官员、社群和家庭可能以不同方式承担征发、保护、迁移或宗教活动；现存来源只能照亮其中一部分，不能替沉默者制造统一的经验。

\eventanchor{EH-0185-SOUTHERN-PALACE-FIRE}{中平二年（185年）·南宫失火，灵帝朝为修复宫室加重财源筹措并扩大官爵交易}账本记录东汉在中平二年（185年）发生“南宫失火，灵帝朝为修复宫室加重财源筹措并扩大官爵交易”。条目把背景概括为黄巾与凉州战事已大量耗费库藏，宫殿火灾又迫使朝廷寻找快速且不稳定的收入，并指出卖官和非常征敛损害官僚任用的合法性，并将财政压力转嫁给地方经营者和求官者。这个节点进入区域社会时，要经由文书、道路、军队、市场、寺院和家庭等不同中介；因此，政治行动的宣布、地方接收和日常生活的改变不能被视作同一时刻完成。

本条依据《后汉书·灵帝纪》；《后汉书·宦者列传》；《后汉纪》卷一百一十，证据提示为“火灾与修营事实可靠；财政收入、卖官范围和因果强度有争议”。这些材料可支持年序、官方决定、战事或特定地点的线索，却难以直接统计所有人的财富、迁徙、信仰和动机。更稳妥的读法是区分诏令与实施、军报与人口、行政称谓与自我认同，并让地方材料的缺环限制解释的范围。

由此可见，该事件不是孤立的年表点，而是资源、信息与权力关系重新配置的一环。不同地区的官员、社群和家庭可能以不同方式承担征发、保护、迁移或宗教活动；现存来源只能照亮其中一部分，不能替沉默者制造统一的经验。

\eventanchor{EH-0186-ZHANG-CHUN-REBELLION}{中平三年（186年）·张纯、张举在幽州起兵并自立名号，东北边郡与乌桓关系趋于动荡}账本记录东汉与幽州叛军在中平三年（186年）发生“张纯、张举在幽州起兵并自立名号，东北边郡与乌桓关系趋于动荡”。条目把背景概括为中央忙于黄巾和凉州，幽州地方军人及边郡网络可利用征发压力发动反叛，并指出东北军政资源更趋地方化，公孙瓒等将领的军事地位上升，汉廷难以恢复直接控制。这个节点进入区域社会时，要经由文书、道路、军队、市场、寺院和家庭等不同中介；因此，政治行动的宣布、地方接收和日常生活的改变不能被视作同一时刻完成。

本条依据《后汉书·灵帝纪》；《后汉书·公孙瓒传》；《后汉书·乌桓鲜卑列传》，证据提示为“叛乱发生可靠；自立称号、乌桓参与度和兵力有争议”。这些材料可支持年序、官方决定、战事或特定地点的线索，却难以直接统计所有人的财富、迁徙、信仰和动机。更稳妥的读法是区分诏令与实施、军报与人口、行政称谓与自我认同，并让地方材料的缺环限制解释的范围。

由此可见，该事件不是孤立的年表点，而是资源、信息与权力关系重新配置的一环。不同地区的官员、社群和家庭可能以不同方式承担征发、保护、迁移或宗教活动；现存来源只能照亮其中一部分，不能替沉默者制造统一的经验。

\eventanchor{EH-0188-ZHOUMU-APPOINTMENTS}{中平五年（188年）·朝廷改任州牧、州刺史并授刘焉等较大军政权限，以应对叛乱和州郡失控}账本记录东汉在中平五年（188年）发生“朝廷改任州牧、州刺史并授刘焉等较大军政权限，以应对叛乱和州郡失控”。条目把背景概括为中央常规刺史监督难以统合跨郡军队和财赋，朝廷需要拥有更强协调权的地方长官，并指出州级军政中心获得持续资源，刘焉等人可建立独立班底，帝国危机转为区域军阀竞争。这个节点进入区域社会时，要经由文书、道路、军队、市场、寺院和家庭等不同中介；因此，政治行动的宣布、地方接收和日常生活的改变不能被视作同一时刻完成。

本条依据《后汉书·灵帝纪》；《后汉书·刘焉传》；《后汉纪》卷一百一十三，证据提示为“任命与制度调整可靠；“州牧制”覆盖范围和权力强度有争议”。这些材料可支持年序、官方决定、战事或特定地点的线索，却难以直接统计所有人的财富、迁徙、信仰和动机。更稳妥的读法是区分诏令与实施、军报与人口、行政称谓与自我认同，并让地方材料的缺环限制解释的范围。

由此可见，该事件不是孤立的年表点，而是资源、信息与权力关系重新配置的一环。不同地区的官员、社群和家庭可能以不同方式承担征发、保护、迁移或宗教活动；现存来源只能照亮其中一部分，不能替沉默者制造统一的经验。

\eventanchor{EH-0188-WESTERN-GARDEN-ARMY}{中平五年（188年）·灵帝置西园八校尉，令宦官蹇硕领兵，试图以宫廷军队制衡外朝将领}账本记录东汉在中平五年（188年）发生“灵帝置西园八校尉，令宦官蹇硕领兵，试图以宫廷军队制衡外朝将领”。条目把背景概括为黄巾后外朝将领和州牧军力膨胀，灵帝与宦官集团试图加强直接受控的京师武力，并指出禁军指挥权更碎片化，何进、蹇硕与宦官的冲突为灵帝死后的宫廷暴力预设条件。这个节点进入区域社会时，要经由文书、道路、军队、市场、寺院和家庭等不同中介；因此，政治行动的宣布、地方接收和日常生活的改变不能被视作同一时刻完成。

本条依据《后汉书·灵帝纪》；《后汉书·何进传》；《后汉纪》卷一百一十三，证据提示为“建置事实可靠；八校尉实际兵额和指挥关系有争议”。这些材料可支持年序、官方决定、战事或特定地点的线索，却难以直接统计所有人的财富、迁徙、信仰和动机。更稳妥的读法是区分诏令与实施、军报与人口、行政称谓与自我认同，并让地方材料的缺环限制解释的范围。

由此可见，该事件不是孤立的年表点，而是资源、信息与权力关系重新配置的一环。不同地区的官员、社群和家庭可能以不同方式承担征发、保护、迁移或宗教活动；现存来源只能照亮其中一部分，不能替沉默者制造统一的经验。

\eventanchor{EH-0189-LINGDI-DEATH}{中平六年四月（189年）·灵帝去世，皇子刘辩即位，是为少帝，何太后临朝，何进掌握外朝军权}账本记录东汉在中平六年四月（189年）发生“灵帝去世，皇子刘辩即位，是为少帝，何太后临朝，何进掌握外朝军权”。条目把背景概括为灵帝留下何后所生少帝与王美人所生陈留王两名皇子，外戚和宦官围绕继统与禁军互不信任，并指出何进以大将军身份成为反宦官集团核心，脆弱的少帝政权迅速演变为宫廷武装对抗。这个节点进入区域社会时，要经由文书、道路、军队、市场、寺院和家庭等不同中介；因此，政治行动的宣布、地方接收和日常生活的改变不能被视作同一时刻完成。

本条依据《后汉书·灵帝纪》；《后汉书·何进传》；《后汉纪》卷一百一十四，证据提示为“继位过程可靠；灵帝遗命和何进军权范围有争议”。这些材料可支持年序、官方决定、战事或特定地点的线索，却难以直接统计所有人的财富、迁徙、信仰和动机。更稳妥的读法是区分诏令与实施、军报与人口、行政称谓与自我认同，并让地方材料的缺环限制解释的范围。

由此可见，该事件不是孤立的年表点，而是资源、信息与权力关系重新配置的一环。不同地区的官员、社群和家庭可能以不同方式承担征发、保护、迁移或宗教活动；现存来源只能照亮其中一部分，不能替沉默者制造统一的经验。

\eventanchor{EH-0189-HE-JIN-MASSACRE}{中平六年八月（189年）·何进谋诛宦官而被杀，袁绍等率兵入宫屠戮宦官，少帝与陈留王一度被劫出京}账本记录东汉在中平六年八月（189年）发生“何进谋诛宦官而被杀，袁绍等率兵入宫屠戮宦官，少帝与陈留王一度被劫出京”。条目把背景概括为何进未能协调太后、宦官和禁军，又以外兵施压宫廷，宦官遂先发制人，并指出京师禁中秩序被摧毁，董卓得以率西凉军进入权力真空，皇帝安全不再由朝廷机构保障。这个节点进入区域社会时，要经由文书、道路、军队、市场、寺院和家庭等不同中介；因此，政治行动的宣布、地方接收和日常生活的改变不能被视作同一时刻完成。

本条依据《后汉书·何进传》；《后汉书·袁绍传》；《后汉纪》卷一百一十四，证据提示为“核心经过可靠；何进召董卓入京的时点、屠杀规模和责任分配有争议”。这些材料可支持年序、官方决定、战事或特定地点的线索，却难以直接统计所有人的财富、迁徙、信仰和动机。更稳妥的读法是区分诏令与实施、军报与人口、行政称谓与自我认同，并让地方材料的缺环限制解释的范围。

由此可见，该事件不是孤立的年表点，而是资源、信息与权力关系重新配置的一环。不同地区的官员、社群和家庭可能以不同方式承担征发、保护、迁移或宗教活动；现存来源只能照亮其中一部分，不能替沉默者制造统一的经验。

\eventanchor{EH-0189-DONG-ZHUO-DEPOSITION}{中平六年九月（189年）·董卓废少帝刘辩，立陈留王刘协为献帝，并毒杀何太后}账本记录东汉与董卓集团在中平六年九月（189年）发生“董卓废少帝刘辩，立陈留王刘协为献帝，并毒杀何太后”。条目把背景概括为宫廷屠杀后禁军与外戚力量瓦解，董卓以西凉军和控制皇帝取得压倒性优势，并指出献帝即位而无自主军权，废立打破皇位不可侵犯的预期，各地州牧和将领获得讨卓的政治理由。这个节点进入区域社会时，要经由文书、道路、军队、市场、寺院和家庭等不同中介；因此，政治行动的宣布、地方接收和日常生活的改变不能被视作同一时刻完成。

本条依据《后汉书·献帝纪》；《后汉书·董卓传》；《后汉纪》卷一百一十四，证据提示为“废立事实可靠；董卓威逼朝臣的具体程序与毒杀细节有争议”。这些材料可支持年序、官方决定、战事或特定地点的线索，却难以直接统计所有人的财富、迁徙、信仰和动机。更稳妥的读法是区分诏令与实施、军报与人口、行政称谓与自我认同，并让地方材料的缺环限制解释的范围。

由此可见，该事件不是孤立的年表点，而是资源、信息与权力关系重新配置的一环。不同地区的官员、社群和家庭可能以不同方式承担征发、保护、迁移或宗教活动；现存来源只能照亮其中一部分，不能替沉默者制造统一的经验。

\eventanchor{EH-0190-ANTI-DONG-COALITION}{初平元年（190年）·袁绍等关东将领以讨董为名结盟，东汉中央与地方军事集团公开分裂}账本记录关东诸州郡与董卓集团在初平元年（190年）发生“袁绍等关东将领以讨董为名结盟，东汉中央与地方军事集团公开分裂”。条目把背景概括为董卓废立和京师暴力激化州郡不满，地方长官已拥有募兵、粮饷和声望网络，并指出讨董联盟未能建立统一指挥，却使各地军事集团以“奉汉”名义独立行动，中央财政与命令体系瓦解。这个节点进入区域社会时，要经由文书、道路、军队、市场、寺院和家庭等不同中介；因此，政治行动的宣布、地方接收和日常生活的改变不能被视作同一时刻完成。

本条依据《后汉书·献帝纪》；《后汉书·袁绍传》；《三国志·武帝纪》，证据提示为“联盟形成可靠；成员实际参战、盟约约束与兵力数字有争议”。这些材料可支持年序、官方决定、战事或特定地点的线索，却难以直接统计所有人的财富、迁徙、信仰和动机。更稳妥的读法是区分诏令与实施、军报与人口、行政称谓与自我认同，并让地方材料的缺环限制解释的范围。

由此可见，该事件不是孤立的年表点，而是资源、信息与权力关系重新配置的一环。不同地区的官员、社群和家庭可能以不同方式承担征发、保护、迁移或宗教活动；现存来源只能照亮其中一部分，不能替沉默者制造统一的经验。

\eventanchor{EH-0190-CHANGAN-TRANSFER}{初平元年（190年）·董卓挟献帝西迁长安，并焚毁洛阳宫室与居民区}账本记录董卓控制下的东汉在初平元年（190年）发生“董卓挟献帝西迁长安，并焚毁洛阳宫室与居民区”。条目把背景概括为关东联军逼近且董卓难以守洛阳，西凉军更依赖关中基地和函谷关防线，并指出献帝朝廷与东部财赋、官僚网络分离；洛阳破坏成为国家崩解的象征，具体毁坏程度须与遗址分期相互校验。这个节点进入区域社会时，要经由文书、道路、军队、市场、寺院和家庭等不同中介；因此，政治行动的宣布、地方接收和日常生活的改变不能被视作同一时刻完成。

本条依据《后汉书·献帝纪》；《后汉书·董卓传》；《三国志·董卓传》；汉魏洛阳城遗址考古报告，证据提示为“迁都事实可靠；焚毁范围、人口损失和考古层位对应有争议”。这些材料可支持年序、官方决定、战事或特定地点的线索，却难以直接统计所有人的财富、迁徙、信仰和动机。更稳妥的读法是区分诏令与实施、军报与人口、行政称谓与自我认同，并让地方材料的缺环限制解释的范围。

由此可见，该事件不是孤立的年表点，而是资源、信息与权力关系重新配置的一环。不同地区的官员、社群和家庭可能以不同方式承担征发、保护、迁移或宗教活动；现存来源只能照亮其中一部分，不能替沉默者制造统一的经验。

\eventanchor{EH-0191-YANGREN-BATTLE}{初平二年（191年）·孙坚在阳人一带击败董卓部将，迫使董卓军进一步收缩}账本记录孙坚军与董卓军在初平二年（191年）发生“孙坚在阳人一带击败董卓部将，迫使董卓军进一步收缩”。条目把背景概括为董卓迁都后仍派军东出牵制诸侯，孙坚获得袁术支持并集中荆州、江东兵员，并指出孙坚展示了地方军队的攻战能力，但联军不能统一利用战果，董卓政权未被迅速推翻。这个节点进入区域社会时，要经由文书、道路、军队、市场、寺院和家庭等不同中介；因此，政治行动的宣布、地方接收和日常生活的改变不能被视作同一时刻完成。

本条依据《三国志·孙破虏讨逆传》；《后汉书·董卓传》；《资治通鉴》卷六十，证据提示为“战斗结果大体可靠；战场地望、将领指挥和追击范围有争议”。这些材料可支持年序、官方决定、战事或特定地点的线索，却难以直接统计所有人的财富、迁徙、信仰和动机。更稳妥的读法是区分诏令与实施、军报与人口、行政称谓与自我认同，并让地方材料的缺环限制解释的范围。

由此可见，该事件不是孤立的年表点，而是资源、信息与权力关系重新配置的一环。不同地区的官员、社群和家庭可能以不同方式承担征发、保护、迁移或宗教活动；现存来源只能照亮其中一部分，不能替沉默者制造统一的经验。

\eventanchor{EH-0191-YUAN-SHAO-JIZHOU}{初平二年（191年）·袁绍取得冀州控制权，以河北户口、粮仓和士族网络建立独立基地}账本记录袁绍集团与韩馥在初平二年（191年）发生“袁绍取得冀州控制权，以河北户口、粮仓和士族网络建立独立基地”。条目把背景概括为讨董联盟内部分裂，袁绍需要稳定粮饷和兵员，冀州富庶且韩馥政治基础薄弱，并指出袁绍成为北方最强军事集团之一，讨董名义让位于各集团争夺州郡资源。这个节点进入区域社会时，要经由文书、道路、军队、市场、寺院和家庭等不同中介；因此，政治行动的宣布、地方接收和日常生活的改变不能被视作同一时刻完成。

本条依据《后汉书·袁绍传》；《三国志·袁绍传》；《英雄记》辑文见裴松之注，证据提示为“取得冀州事实可靠；韩馥让州是否出于自愿和具体交易有争议”。这些材料可支持年序、官方决定、战事或特定地点的线索，却难以直接统计所有人的财富、迁徙、信仰和动机。更稳妥的读法是区分诏令与实施、军报与人口、行政称谓与自我认同，并让地方材料的缺环限制解释的范围。

由此可见，该事件不是孤立的年表点，而是资源、信息与权力关系重新配置的一环。不同地区的官员、社群和家庭可能以不同方式承担征发、保护、迁移或宗教活动；现存来源只能照亮其中一部分，不能替沉默者制造统一的经验。

\eventanchor{EH-0192-DONG-ZHUO-ASSASSINATION}{初平三年四月（192年）·王允、吕布合谋杀董卓，长安朝廷短暂摆脱董卓直接控制}账本记录长安朝廷与董卓集团在初平三年四月（192年）发生“王允、吕布合谋杀董卓，长安朝廷短暂摆脱董卓直接控制”。条目把背景概括为董卓专权、军纪和迁都损害朝廷与关中社会，王允利用吕布与董卓的矛盾发动宫廷袭击，并指出董卓死后其西凉旧部并未解体，王允缺乏整合军队的能力，长安政权更易遭反攻。这个节点进入区域社会时，要经由文书、道路、军队、市场、寺院和家庭等不同中介；因此，政治行动的宣布、地方接收和日常生活的改变不能被视作同一时刻完成。

本条依据《后汉书·董卓传》；《三国志·吕布传》；《后汉纪》卷一百一十七，证据提示为“刺杀日期和结果可靠；王允、吕布动机及貂蝉故事有争议与后世加工”。这些材料可支持年序、官方决定、战事或特定地点的线索，却难以直接统计所有人的财富、迁徙、信仰和动机。更稳妥的读法是区分诏令与实施、军报与人口、行政称谓与自我认同，并让地方材料的缺环限制解释的范围。

由此可见，该事件不是孤立的年表点，而是资源、信息与权力关系重新配置的一环。不同地区的官员、社群和家庭可能以不同方式承担征发、保护、迁移或宗教活动；现存来源只能照亮其中一部分，不能替沉默者制造统一的经验。

\eventanchor{EH-0192-LI-JUE-CHANGAN}{初平三年（192年）·李傕、郭汜等董卓旧部攻入长安，杀王允并控制献帝}账本记录李傕郭汜集团与长安朝廷在初平三年（192年）发生“李傕、郭汜等董卓旧部攻入长安，杀王允并控制献帝”。条目把背景概括为王允拒绝或未能妥善安置董卓旧部，西凉军以生存和报复为由重新聚集，并指出献帝再度成为军人集团手中的政治资源，长安朝廷难以向关东发出有效命令。这个节点进入区域社会时，要经由文书、道路、军队、市场、寺院和家庭等不同中介；因此，政治行动的宣布、地方接收和日常生活的改变不能被视作同一时刻完成。

本条依据《后汉书·献帝纪》；《后汉书·董卓传》；《三国志·董卓传》，证据提示为“攻陷长安事实可靠；兵力、围城细节和各将责任有争议”。这些材料可支持年序、官方决定、战事或特定地点的线索，却难以直接统计所有人的财富、迁徙、信仰和动机。更稳妥的读法是区分诏令与实施、军报与人口、行政称谓与自我认同，并让地方材料的缺环限制解释的范围。

由此可见，该事件不是孤立的年表点，而是资源、信息与权力关系重新配置的一环。不同地区的官员、社群和家庭可能以不同方式承担征发、保护、迁移或宗教活动；现存来源只能照亮其中一部分，不能替沉默者制造统一的经验。

\eventanchor{EH-0193-CAO-CAO-XUZHOU}{初平四年（193年）·曹操以父曹嵩遇害为由两次攻徐州，沿途杀掠记载引发史料评价争议}账本记录曹操集团与徐州陶谦在初平四年（193年）发生“曹操以父曹嵩遇害为由两次攻徐州，沿途杀掠记载引发史料评价争议”。条目把背景概括为曹嵩死于徐州境内激化曹操报复，曹操亦需夺取东部粮赋和交通资源，并指出徐州社会遭受严重破坏，曹操因后方空虚未能吞并徐州，吕布得以进入兖州。这个节点进入区域社会时，要经由文书、道路、军队、市场、寺院和家庭等不同中介；因此，政治行动的宣布、地方接收和日常生活的改变不能被视作同一时刻完成。

本条依据《三国志·武帝纪》；《后汉书·陶谦传》；《资治通鉴》卷六十一，证据提示为“出兵缘由与战事可靠；屠杀范围、死亡数字和陶谦责任有争议”。这些材料可支持年序、官方决定、战事或特定地点的线索，却难以直接统计所有人的财富、迁徙、信仰和动机。更稳妥的读法是区分诏令与实施、军报与人口、行政称谓与自我认同，并让地方材料的缺环限制解释的范围。

由此可见，该事件不是孤立的年表点，而是资源、信息与权力关系重新配置的一环。不同地区的官员、社群和家庭可能以不同方式承担征发、保护、迁移或宗教活动；现存来源只能照亮其中一部分，不能替沉默者制造统一的经验。

\eventanchor{EH-0194-LU-BU-YANZHOU}{兴平元年（194年）·吕布在陈宫等响应下进入兖州，占据曹操后方多处郡县}账本记录吕布、陈宫与曹操集团在兴平元年（194年）发生“吕布在陈宫等响应下进入兖州，占据曹操后方多处郡县”。条目把背景概括为曹操主力在徐州作战，兖州士人与地方官对其征发和战事负担不满，吕布获得可乘之机，并指出曹操被迫回师争夺根据地，兖州控制权成为其能否继续扩张的决定性资源。这个节点进入区域社会时，要经由文书、道路、军队、市场、寺院和家庭等不同中介；因此，政治行动的宣布、地方接收和日常生活的改变不能被视作同一时刻完成。

本条依据《三国志·吕布传》；《三国志·武帝纪》；《后汉书·吕布传》，证据提示为“吕布入兖州可靠；响应郡县、陈宫作用和地方支持程度有争议”。这些材料可支持年序、官方决定、战事或特定地点的线索，却难以直接统计所有人的财富、迁徙、信仰和动机。更稳妥的读法是区分诏令与实施、军报与人口、行政称谓与自我认同，并让地方材料的缺环限制解释的范围。

由此可见，该事件不是孤立的年表点，而是资源、信息与权力关系重新配置的一环。不同地区的官员、社群和家庭可能以不同方式承担征发、保护、迁移或宗教活动；现存来源只能照亮其中一部分，不能替沉默者制造统一的经验。

\eventanchor{EH-0195-LI-GUO-CONFLICT}{兴平二年（195年）·李傕、郭汜相互攻战并挟持献帝，长安政权陷入内战}账本记录李傕郭汜集团与献帝朝廷在兴平二年（195年）发生“李傕、郭汜相互攻战并挟持献帝，长安政权陷入内战”。条目把背景概括为西凉军集团缺少共享的财政和指挥机制，个人猜忌使对献帝的共同控制转化为竞争，并指出献帝得以东逃但朝廷威信进一步破产，关东军阀争夺“奉天子”名义的条件成熟。这个节点进入区域社会时，要经由文书、道路、军队、市场、寺院和家庭等不同中介；因此，政治行动的宣布、地方接收和日常生活的改变不能被视作同一时刻完成。

本条依据《后汉书·献帝纪》；《三国志·董卓传》；《后汉纪》卷一百二十，证据提示为“内战事实可靠；冲突起因、质子细节和各部人数有争议”。这些材料可支持年序、官方决定、战事或特定地点的线索，却难以直接统计所有人的财富、迁徙、信仰和动机。更稳妥的读法是区分诏令与实施、军报与人口、行政称谓与自我认同，并让地方材料的缺环限制解释的范围。

由此可见，该事件不是孤立的年表点，而是资源、信息与权力关系重新配置的一环。不同地区的官员、社群和家庭可能以不同方式承担征发、保护、迁移或宗教活动；现存来源只能照亮其中一部分，不能替沉默者制造统一的经验。

\eventanchor{EH-0196-CAO-CAO-RECEIVES-EMPEROR}{建安元年（196年）·曹操迎献帝至许县，随后以朝廷名义发布诏令并组织屯田}账本记录献帝朝廷与曹操集团在建安元年（196年）发生“曹操迎献帝至许县，随后以朝廷名义发布诏令并组织屯田”。条目把背景概括为献帝东归后缺乏稳定粮饷和保护者，曹操控制兖州后具备粮仓、军队和行政人员，并指出曹操取得诏令合法性和中央名义，许都屯田与军政组织为其北方竞争提供较持续的资源。这个节点进入区域社会时，要经由文书、道路、军队、市场、寺院和家庭等不同中介；因此，政治行动的宣布、地方接收和日常生活的改变不能被视作同一时刻完成。

本条依据《后汉书·献帝纪》；《三国志·武帝纪》；《三国志·任峻传》，证据提示为“迎驾与迁许事实可靠；屯田初创时间、规模和曹操个人作用有争议”。这些材料可支持年序、官方决定、战事或特定地点的线索，却难以直接统计所有人的财富、迁徙、信仰和动机。更稳妥的读法是区分诏令与实施、军报与人口、行政称谓与自我认同，并让地方材料的缺环限制解释的范围。

由此可见，该事件不是孤立的年表点，而是资源、信息与权力关系重新配置的一环。不同地区的官员、社群和家庭可能以不同方式承担征发、保护、迁移或宗教活动；现存来源只能照亮其中一部分，不能替沉默者制造统一的经验。

\eventanchor{EH-0197-YUAN-SHU-EMPEROR}{建安二年（197年）·袁术在寿春称帝，国号仲氏，遭各方军事与外交孤立}账本记录袁术政权与汉廷诸集团在建安二年（197年）发生“袁术在寿春称帝，国号仲氏，遭各方军事与外交孤立”。条目把背景概括为袁术拥有淮南资源并以传国玺、谶纬为号召，误判其他军阀会接受新的皇帝中心，并指出称帝破坏其与孙策、吕布等的合作，反对者可借奉汉名义联合，仲氏迅速失去支持。这个节点进入区域社会时，要经由文书、道路、军队、市场、寺院和家庭等不同中介；因此，政治行动的宣布、地方接收和日常生活的改变不能被视作同一时刻完成。

本条依据《后汉书·袁术传》；《三国志·袁术传》；《资治通鉴》卷六十二，证据提示为“称帝年份可靠；受命符瑞、国号含义和实际统治范围有争议”。这些材料可支持年序、官方决定、战事或特定地点的线索，却难以直接统计所有人的财富、迁徙、信仰和动机。更稳妥的读法是区分诏令与实施、军报与人口、行政称谓与自我认同，并让地方材料的缺环限制解释的范围。

由此可见，该事件不是孤立的年表点，而是资源、信息与权力关系重新配置的一环。不同地区的官员、社群和家庭可能以不同方式承担征发、保护、迁移或宗教活动；现存来源只能照亮其中一部分，不能替沉默者制造统一的经验。

\eventanchor{EH-0198-XIAPI-FALL}{建安三年（198年）·曹操与刘备围下邳，吕布被部下缚降后处死}账本记录曹操、刘备与吕布集团在建安三年（198年）发生“曹操与刘备围下邳，吕布被部下缚降后处死”。条目把背景概括为吕布在徐州与刘备、袁术反复结盟，粮饷和部将忠诚均不稳定，曹操得以联合刘备围攻，并指出曹操消除东部强敌并控制徐州要地，为集中兵力同袁绍争夺北方创造条件。这个节点进入区域社会时，要经由文书、道路、军队、市场、寺院和家庭等不同中介；因此，政治行动的宣布、地方接收和日常生活的改变不能被视作同一时刻完成。

本条依据《三国志·吕布传》；《三国志·武帝纪》；《后汉书·吕布传》，证据提示为“战役结果可靠；水攻过程、降将动机和处决细节有争议”。这些材料可支持年序、官方决定、战事或特定地点的线索，却难以直接统计所有人的财富、迁徙、信仰和动机。更稳妥的读法是区分诏令与实施、军报与人口、行政称谓与自我认同，并让地方材料的缺环限制解释的范围。

由此可见，该事件不是孤立的年表点，而是资源、信息与权力关系重新配置的一环。不同地区的官员、社群和家庭可能以不同方式承担征发、保护、迁移或宗教活动；现存来源只能照亮其中一部分，不能替沉默者制造统一的经验。

\eventanchor{EH-0200-GUANDU}{建安五年（200年）·曹操在官渡击败袁绍主力，焚其乌巢粮仓并扭转兵力劣势}账本记录曹操集团与袁绍集团在建安五年（200年）发生“曹操在官渡击败袁绍主力，焚其乌巢粮仓并扭转兵力劣势”。条目把背景概括为袁绍据河北人口粮赋南下，曹操依许都朝廷和屯田维持防线，但面临补给和兵力压力，并指出袁绍主力溃散，曹操获得逐步吞并河北的战略窗口；史书兵数多具夸张与政治修辞。这个节点进入区域社会时，要经由文书、道路、军队、市场、寺院和家庭等不同中介；因此，政治行动的宣布、地方接收和日常生活的改变不能被视作同一时刻完成。

本条依据《三国志·武帝纪》；《三国志·袁绍传》；《后汉书·袁绍传》，证据提示为“战役年份和胜负可靠；双方兵数、许攸来降与乌巢细节有争议”。这些材料可支持年序、官方决定、战事或特定地点的线索，却难以直接统计所有人的财富、迁徙、信仰和动机。更稳妥的读法是区分诏令与实施、军报与人口、行政称谓与自我认同，并让地方材料的缺环限制解释的范围。

由此可见，该事件不是孤立的年表点，而是资源、信息与权力关系重新配置的一环。不同地区的官员、社群和家庭可能以不同方式承担征发、保护、迁移或宗教活动；现存来源只能照亮其中一部分，不能替沉默者制造统一的经验。

\eventanchor{EH-0202-YUAN-SHAO-DEATH}{建安七年（202年）·袁绍去世，袁谭、袁尚围绕继承与河北资源发生内斗}账本记录袁绍集团在建安七年（202年）发生“袁绍去世，袁谭、袁尚围绕继承与河北资源发生内斗”。条目把背景概括为官渡失败削弱袁绍威望，且其诸子、审配和逢纪等集团对继承缺少共同安排，并指出河北军政资源被内斗分割，曹操可利用降将与军事压力逐步夺取黄河以北。这个节点进入区域社会时，要经由文书、道路、军队、市场、寺院和家庭等不同中介；因此，政治行动的宣布、地方接收和日常生活的改变不能被视作同一时刻完成。

本条依据《三国志·袁绍传》；《后汉书·袁绍传》；《资治通鉴》卷六十四，证据提示为“卒年可靠；继承遗命真伪和兄弟冲突责任有争议”。这些材料可支持年序、官方决定、战事或特定地点的线索，却难以直接统计所有人的财富、迁徙、信仰和动机。更稳妥的读法是区分诏令与实施、军报与人口、行政称谓与自我认同，并让地方材料的缺环限制解释的范围。

由此可见，该事件不是孤立的年表点，而是资源、信息与权力关系重新配置的一环。不同地区的官员、社群和家庭可能以不同方式承担征发、保护、迁移或宗教活动；现存来源只能照亮其中一部分，不能替沉默者制造统一的经验。

\eventanchor{EH-0204-TAKE-YE}{建安九年（204年）·曹操攻克邺城，取得冀州治所、仓储与袁绍集团的核心行政资源}账本记录曹操集团与袁尚集团在建安九年（204年）发生“曹操攻克邺城，取得冀州治所、仓储与袁绍集团的核心行政资源”。条目把背景概括为袁氏兄弟争斗削弱协防，曹操以官渡后积累的军政资源围攻冀州中心，并指出曹操获得河北粮仓、人口和文士网络，邺城后来成为其北方政权的重要基地。这个节点进入区域社会时，要经由文书、道路、军队、市场、寺院和家庭等不同中介；因此，政治行动的宣布、地方接收和日常生活的改变不能被视作同一时刻完成。

本条依据《三国志·武帝纪》；《三国志·袁绍传》；《后汉书·袁绍传》，证据提示为“攻城结果可靠；围城时长、灾情和降附过程有争议”。这些材料可支持年序、官方决定、战事或特定地点的线索，却难以直接统计所有人的财富、迁徙、信仰和动机。更稳妥的读法是区分诏令与实施、军报与人口、行政称谓与自我认同，并让地方材料的缺环限制解释的范围。

由此可见，该事件不是孤立的年表点，而是资源、信息与权力关系重新配置的一环。不同地区的官员、社群和家庭可能以不同方式承担征发、保护、迁移或宗教活动；现存来源只能照亮其中一部分，不能替沉默者制造统一的经验。

\eventanchor{EH-0205-YUAN-TAN-DEFEAT}{建安十年（205年）·曹操在南皮击败并杀袁谭，袁氏在冀州的主要抵抗瓦解}账本记录曹操集团与袁谭集团在建安十年（205年）发生“曹操在南皮击败并杀袁谭，袁氏在冀州的主要抵抗瓦解”。条目把背景概括为袁谭为对抗袁尚曾借助曹操，随后又与曹操争夺冀州，失去稳定联盟和补给，并指出曹操进一步统一河北平原，袁尚、袁熙只能退向幽州和乌桓，北方统一趋势加强。这个节点进入区域社会时，要经由文书、道路、军队、市场、寺院和家庭等不同中介；因此，政治行动的宣布、地方接收和日常生活的改变不能被视作同一时刻完成。

本条依据《三国志·武帝纪》；《三国志·袁绍传》；《后汉书·袁绍传》，证据提示为“战役结果可靠；袁谭求援与战死细节有争议”。这些材料可支持年序、官方决定、战事或特定地点的线索，却难以直接统计所有人的财富、迁徙、信仰和动机。更稳妥的读法是区分诏令与实施、军报与人口、行政称谓与自我认同，并让地方材料的缺环限制解释的范围。

由此可见，该事件不是孤立的年表点，而是资源、信息与权力关系重新配置的一环。不同地区的官员、社群和家庭可能以不同方式承担征发、保护、迁移或宗教活动；现存来源只能照亮其中一部分，不能替沉默者制造统一的经验。

\eventanchor{EH-0207-WUHUAN-CAMPAIGN}{建安十二年（207年）·曹操北征乌桓，在白狼山击破蹋顿，袁尚、袁熙残部失去依托}账本记录曹操集团、乌桓与袁氏残部在建安十二年（207年）发生“曹操北征乌桓，在白狼山击破蹋顿，袁尚、袁熙残部失去依托”。条目把背景概括为袁氏残部依托乌桓威胁幽州，曹操需先解除北方侧翼才可南下争夺荆州，并指出乌桓联盟受挫，曹操基本控制河北与幽州通道；边疆部众的后续安置不能简化为完全臣服。这个节点进入区域社会时，要经由文书、道路、军队、市场、寺院和家庭等不同中介；因此，政治行动的宣布、地方接收和日常生活的改变不能被视作同一时刻完成。

本条依据《三国志·武帝纪》；《三国志·乌丸鲜卑东夷传》；《后汉书·乌桓鲜卑列传》，证据提示为“北征和胜负可靠；行军路线、参战部众和白狼山地望有争议”。这些材料可支持年序、官方决定、战事或特定地点的线索，却难以直接统计所有人的财富、迁徙、信仰和动机。更稳妥的读法是区分诏令与实施、军报与人口、行政称谓与自我认同，并让地方材料的缺环限制解释的范围。

由此可见，该事件不是孤立的年表点，而是资源、信息与权力关系重新配置的一环。不同地区的官员、社群和家庭可能以不同方式承担征发、保护、迁移或宗教活动；现存来源只能照亮其中一部分，不能替沉默者制造统一的经验。

\subsection{边疆、战争与移动}

\eventanchor{EH-0208-JINGZHOU-TRANSFER}{建安十三年（208年）·刘表去世后，刘琮降曹，曹操取得荆州北部并南下追击刘备}账本记录曹操集团与荆州刘表集团在建安十三年（208年）发生“刘表去世后，刘琮降曹，曹操取得荆州北部并南下追击刘备”。条目把背景概括为刘表长期维持荆州相对独立，病逝后继承集团分裂，曹操以北方优势迅速压迫襄阳，并指出曹操获得部分荆州军政资源并逼近长江，孙权与刘备面临结盟或各自被吞并的选择。这个节点进入区域社会时，要经由文书、道路、军队、市场、寺院和家庭等不同中介；因此，政治行动的宣布、地方接收和日常生活的改变不能被视作同一时刻完成。

本条依据《三国志·武帝纪》；《三国志·先主传》；《后汉书·刘表传》，证据提示为“刘表卒与刘琮降曹可靠；降附动机、蔡氏作用和荆州实际交接范围有争议”。这些材料可支持年序、官方决定、战事或特定地点的线索，却难以直接统计所有人的财富、迁徙、信仰和动机。更稳妥的读法是区分诏令与实施、军报与人口、行政称谓与自我认同，并让地方材料的缺环限制解释的范围。

由此可见，该事件不是孤立的年表点，而是资源、信息与权力关系重新配置的一环。不同地区的官员、社群和家庭可能以不同方式承担征发、保护、迁移或宗教活动；现存来源只能照亮其中一部分，不能替沉默者制造统一的经验。

\eventanchor{EH-0208-RED-CLIFFS}{建安十三年冬（208年）·孙刘联军在赤壁一带击败曹操水军，曹操撤回北方}账本记录曹操集团与孙权刘备联军在建安十三年冬（208年）发生“孙刘联军在赤壁一带击败曹操水军，曹操撤回北方”。条目把背景概括为曹操快速南下后水军、疫病与江淮—荆州水域经验不足，孙刘联盟则须阻止其控制长江，并指出曹操未能一举统一南方，孙权和刘备各自取得扩张空间；演义化细节不可替代纪传材料。这个节点进入区域社会时，要经由文书、道路、军队、市场、寺院和家庭等不同中介；因此，政治行动的宣布、地方接收和日常生活的改变不能被视作同一时刻完成。

本条依据《三国志·武帝纪》；《三国志·周瑜传》；《三国志·先主传》；《资治通鉴》卷六十五，证据提示为“战役结果可靠；战场位置、火攻执行者、疫病影响和兵数有争议”。这些材料可支持年序、官方决定、战事或特定地点的线索，却难以直接统计所有人的财富、迁徙、信仰和动机。更稳妥的读法是区分诏令与实施、军报与人口、行政称谓与自我认同，并让地方材料的缺环限制解释的范围。

由此可见，该事件不是孤立的年表点，而是资源、信息与权力关系重新配置的一环。不同地区的官员、社群和家庭可能以不同方式承担征发、保护、迁移或宗教活动；现存来源只能照亮其中一部分，不能替沉默者制造统一的经验。

\eventanchor{EH-0209-JINGZHOU-DIVISION}{建安十四年（209年）·孙刘在荆州南部及江陵周边重新分配城郡与防线，三方沿长江形成对峙}账本记录孙权、刘备与曹操集团在建安十四年（209年）发生“孙刘在荆州南部及江陵周边重新分配城郡与防线，三方沿长江形成对峙”。条目把背景概括为赤壁后曹操退守北方，孙刘需在共同抗曹与各自获得荆州人口、港口和粮赋之间协商，并指出荆州成为孙刘联盟的利益冲突核心，三方边界并非固定，后续战争围绕其交通和郡县展开。这个节点进入区域社会时，要经由文书、道路、军队、市场、寺院和家庭等不同中介；因此，政治行动的宣布、地方接收和日常生活的改变不能被视作同一时刻完成。

本条依据《三国志·周瑜传》；《三国志·先主传》；《三国志·鲁肃传》，证据提示为“控制格局大体可靠；“借荆州”条款、具体郡界和交接时间有争议”。这些材料可支持年序、官方决定、战事或特定地点的线索，却难以直接统计所有人的财富、迁徙、信仰和动机。更稳妥的读法是区分诏令与实施、军报与人口、行政称谓与自我认同，并让地方材料的缺环限制解释的范围。

由此可见，该事件不是孤立的年表点，而是资源、信息与权力关系重新配置的一环。不同地区的官员、社群和家庭可能以不同方式承担征发、保护、迁移或宗教活动；现存来源只能照亮其中一部分，不能替沉默者制造统一的经验。

\eventanchor{EH-0211-TONG-PASS}{建安十六年（211年）·曹操在潼关击败马超、韩遂等关中联军，关中军事集团分裂}账本记录曹操集团与关中诸将在建安十六年（211年）发生“曹操在潼关击败马超、韩遂等关中联军，关中军事集团分裂”。条目把背景概括为曹操西进威胁关中诸将的地盘与补给，马超、韩遂等暂时联合以守关中门户，并指出关中联军瓦解，曹操取得西向战略主动，但凉州与汉中仍非完全受控区域。这个节点进入区域社会时，要经由文书、道路、军队、市场、寺院和家庭等不同中介；因此，政治行动的宣布、地方接收和日常生活的改变不能被视作同一时刻完成。

本条依据《三国志·武帝纪》；《三国志·马超传》；《后汉书·董卓传》，证据提示为“战役结果可靠；离间计、渡河路线和联军构成有争议”。这些材料可支持年序、官方决定、战事或特定地点的线索，却难以直接统计所有人的财富、迁徙、信仰和动机。更稳妥的读法是区分诏令与实施、军报与人口、行政称谓与自我认同，并让地方材料的缺环限制解释的范围。

由此可见，该事件不是孤立的年表点，而是资源、信息与权力关系重新配置的一环。不同地区的官员、社群和家庭可能以不同方式承担征发、保护、迁移或宗教活动；现存来源只能照亮其中一部分，不能替沉默者制造统一的经验。

\eventanchor{EH-0212-LIU-BEI-YIZHOU}{建安十七年（212年）·刘备由葭萌进军益州，与刘璋决裂，益州内战开始}账本记录刘备集团与益州刘璋政权在建安十七年（212年）发生“刘备由葭萌进军益州，与刘璋决裂，益州内战开始”。条目把背景概括为刘璋为防张鲁邀刘备入蜀，却无法协调益州旧有官僚、豪强和外来刘备军的利益，并指出刘备获得争夺巴蜀的立足点，益州人口、粮赋与山地关隘成为三方格局的关键资源。这个节点进入区域社会时，要经由文书、道路、军队、市场、寺院和家庭等不同中介；因此，政治行动的宣布、地方接收和日常生活的改变不能被视作同一时刻完成。

本条依据《三国志·先主传》；《三国志·刘璋传》；《华阳国志·刘先主志》，证据提示为“进军与决裂事实可靠；刘备受邀初衷、庞统建议和开战责任有争议”。这些材料可支持年序、官方决定、战事或特定地点的线索，却难以直接统计所有人的财富、迁徙、信仰和动机。更稳妥的读法是区分诏令与实施、军报与人口、行政称谓与自我认同，并让地方材料的缺环限制解释的范围。

由此可见，该事件不是孤立的年表点，而是资源、信息与权力关系重新配置的一环。不同地区的官员、社群和家庭可能以不同方式承担征发、保护、迁移或宗教活动；现存来源只能照亮其中一部分，不能替沉默者制造统一的经验。

\eventanchor{EH-0214-CHENGDU-SURRENDER}{建安十九年（214年）·刘璋降刘备，刘备据成都并接收益州行政与军需资源}账本记录刘备集团与益州刘璋政权在建安十九年（214年）发生“刘璋降刘备，刘备据成都并接收益州行政与军需资源”。条目把背景概括为刘备军控制益州要地并获部分地方官民支持，刘璋难以组织持续反攻或外援，并指出刘备获得稳固的粮赋与兵员基地，可与曹操、孙权并列竞争；旧益州集团的整合仍须通过任官和联盟完成。这个节点进入区域社会时，要经由文书、道路、军队、市场、寺院和家庭等不同中介；因此，政治行动的宣布、地方接收和日常生活的改变不能被视作同一时刻完成。

本条依据《三国志·先主传》；《三国志·刘璋传》；《华阳国志·刘先主志》，证据提示为“降附年份可靠；围城谈判、地方士族选择和接收程度有争议”。这些材料可支持年序、官方决定、战事或特定地点的线索，却难以直接统计所有人的财富、迁徙、信仰和动机。更稳妥的读法是区分诏令与实施、军报与人口、行政称谓与自我认同，并让地方材料的缺环限制解释的范围。

由此可见，该事件不是孤立的年表点，而是资源、信息与权力关系重新配置的一环。不同地区的官员、社群和家庭可能以不同方式承担征发、保护、迁移或宗教活动；现存来源只能照亮其中一部分，不能替沉默者制造统一的经验。

\eventanchor{EH-0215-CAO-CAO-HANZHONG}{建安二十年（215年）·曹操攻取汉中，张鲁降，曹操控制连接关中与巴蜀的北部通道}账本记录曹操集团与张鲁政权在建安二十年（215年）发生“曹操攻取汉中，张鲁降，曹操控制连接关中与巴蜀的北部通道”。条目把背景概括为张鲁政权控制汉中多年，曹操在平定关中后需消除其对关中南缘与益州北门的威胁，并指出曹操获得汉中并迁徙部分居民，刘备感到益州北部受压，汉中成为双方下一轮战争焦点。这个节点进入区域社会时，要经由文书、道路、军队、市场、寺院和家庭等不同中介；因此，政治行动的宣布、地方接收和日常生活的改变不能被视作同一时刻完成。

本条依据《三国志·武帝纪》；《三国志·张鲁传》；《后汉书·张鲁传》，证据提示为“攻取与降附事实可靠；张鲁宗教组织、人口迁徙和曹操留守部署有争议”。这些材料可支持年序、官方决定、战事或特定地点的线索，却难以直接统计所有人的财富、迁徙、信仰和动机。更稳妥的读法是区分诏令与实施、军报与人口、行政称谓与自我认同，并让地方材料的缺环限制解释的范围。

由此可见，该事件不是孤立的年表点，而是资源、信息与权力关系重新配置的一环。不同地区的官员、社群和家庭可能以不同方式承担征发、保护、迁移或宗教活动；现存来源只能照亮其中一部分，不能替沉默者制造统一的经验。

\eventanchor{EH-0215-SUN-QUAN-JINGZHOU}{建安二十年（215年）·孙权夺取长沙、零陵、桂阳等荆南郡，孙刘以湘水附近重新划界}账本记录孙权与刘备集团在建安二十年（215年）发生“孙权夺取长沙、零陵、桂阳等荆南郡，孙刘以湘水附近重新划界”。条目把背景概括为刘备据益州后仍控制部分荆州，孙权认为赤壁后分配失衡并需保障长江上游安全，并指出孙刘联盟虽暂时修复以共同对曹，却将荆州主权冲突制度化，关羽日后北进缺少稳固后方。这个节点进入区域社会时，要经由文书、道路、军队、市场、寺院和家庭等不同中介；因此，政治行动的宣布、地方接收和日常生活的改变不能被视作同一时刻完成。

本条依据《三国志·先主传》；《三国志·孙权传》；《三国志·鲁肃传》，证据提示为“夺郡与和议大体可靠；各郡先后、湘水划界位置和实际执行有争议”。这些材料可支持年序、官方决定、战事或特定地点的线索，却难以直接统计所有人的财富、迁徙、信仰和动机。更稳妥的读法是区分诏令与实施、军报与人口、行政称谓与自我认同，并让地方材料的缺环限制解释的范围。

由此可见，该事件不是孤立的年表点，而是资源、信息与权力关系重新配置的一环。不同地区的官员、社群和家庭可能以不同方式承担征发、保护、迁移或宗教活动；现存来源只能照亮其中一部分，不能替沉默者制造统一的经验。

\eventanchor{EH-0216-CAO-CAO-KING-WEI}{建安二十一年（216年）·献帝封曹操为魏王，加九锡，魏国官署与封国体制进一步完备}账本记录献帝朝廷与曹操集团在建安二十一年（216年）发生“献帝封曹操为魏王，加九锡，魏国官署与封国体制进一步完备”。条目把背景概括为曹操已控制北方、朝廷和主要军政资源，需以封国名义安置功臣并提升自身合法性，并指出魏国成为可替代汉朝廷的政治框架，汉帝保留名义而实际权力继续向曹氏集中。这个节点进入区域社会时，要经由文书、道路、军队、市场、寺院和家庭等不同中介；因此，政治行动的宣布、地方接收和日常生活的改变不能被视作同一时刻完成。

本条依据《后汉书·献帝纪》；《三国志·武帝纪》；《三国志·荀彧传》，证据提示为“封王年月可靠；九锡礼仪、魏国机构权能和献帝主动性有争议”。这些材料可支持年序、官方决定、战事或特定地点的线索，却难以直接统计所有人的财富、迁徙、信仰和动机。更稳妥的读法是区分诏令与实施、军报与人口、行政称谓与自我认同，并让地方材料的缺环限制解释的范围。

由此可见，该事件不是孤立的年表点，而是资源、信息与权力关系重新配置的一环。不同地区的官员、社群和家庭可能以不同方式承担征发、保护、迁移或宗教活动；现存来源只能照亮其中一部分，不能替沉默者制造统一的经验。

\eventanchor{EH-0217-JIANYAN-PLAGUE}{建安二十二年（217年）·中原发生大疫，曹丕、王粲等文士和军民遭受损失，战争动员承受额外压力}账本记录曹操控制区在建安二十二年（217年）发生“中原发生大疫，曹丕、王粲等文士和军民遭受损失，战争动员承受额外压力”。条目把背景概括为连年征战、人口流动、营垒与灾害使疾病传播风险上升，官府救济与医疗记录有限，并指出文士群体的死亡留下可见文本记录，军政与地方社会的损失难以量化，不能以文学哀悼推算全国死亡率。这个节点进入区域社会时，要经由文书、道路、军队、市场、寺院和家庭等不同中介；因此，政治行动的宣布、地方接收和日常生活的改变不能被视作同一时刻完成。

本条依据《三国志·武帝纪》；曹植《说疫气》；《文选》所收建安文人作品，证据提示为“疫病流行可靠；死亡规模、病原和区域范围无法确定”。这些材料可支持年序、官方决定、战事或特定地点的线索，却难以直接统计所有人的财富、迁徙、信仰和动机。更稳妥的读法是区分诏令与实施、军报与人口、行政称谓与自我认同，并让地方材料的缺环限制解释的范围。

由此可见，该事件不是孤立的年表点，而是资源、信息与权力关系重新配置的一环。不同地区的官员、社群和家庭可能以不同方式承担征发、保护、迁移或宗教活动；现存来源只能照亮其中一部分，不能替沉默者制造统一的经验。

\eventanchor{EH-0218-HANZHONG-CAMPAIGN}{建安二十三年（218年）·刘备在汉中持续进攻夏侯渊、张郃等部，争夺定军山和阳平关方向的通道}账本记录刘备集团与曹操集团在建安二十三年（218年）发生“刘备在汉中持续进攻夏侯渊、张郃等部，争夺定军山和阳平关方向的通道”。条目把背景概括为曹操占汉中后威胁益州北部，刘备需保住巴蜀门户并取得可防御的山地屏障，并指出汉中战争消耗双方兵力和粮运，为次年刘备取得汉中及关羽北伐创造相互牵连的战场条件。这个节点进入区域社会时，要经由文书、道路、军队、市场、寺院和家庭等不同中介；因此，政治行动的宣布、地方接收和日常生活的改变不能被视作同一时刻完成。

本条依据《三国志·先主传》；《三国志·夏侯渊传》；《三国志·张郃传》，证据提示为“战役阶段可靠；行军路线、各将兵额和地方协力程度有争议”。这些材料可支持年序、官方决定、战事或特定地点的线索，却难以直接统计所有人的财富、迁徙、信仰和动机。更稳妥的读法是区分诏令与实施、军报与人口、行政称谓与自我认同，并让地方材料的缺环限制解释的范围。

由此可见，该事件不是孤立的年表点，而是资源、信息与权力关系重新配置的一环。不同地区的官员、社群和家庭可能以不同方式承担征发、保护、迁移或宗教活动；现存来源只能照亮其中一部分，不能替沉默者制造统一的经验。

\eventanchor{EH-0219-DINGJUN-HANZHONG}{建安二十四年（219年）·刘备军在定军山击杀夏侯渊，曹操亲至后仍撤出汉中，刘备据汉中称汉中王}账本记录刘备集团与曹操集团在建安二十四年（219年）发生“刘备军在定军山击杀夏侯渊，曹操亲至后仍撤出汉中，刘备据汉中称汉中王”。条目把背景概括为刘备长期依托益州供给围攻汉中，夏侯渊前出部署使曹军失去关键统帅，并指出刘备获得汉中防线与“汉中王”名号，三方竞争从州郡割据进一步转向准王朝建构。这个节点进入区域社会时，要经由文书、道路、军队、市场、寺院和家庭等不同中介；因此，政治行动的宣布、地方接收和日常生活的改变不能被视作同一时刻完成。

本条依据《三国志·先主传》；《三国志·夏侯渊传》；《三国志·法正传》，证据提示为“核心结果可靠；黄忠战功、战场细节和称王礼仪有争议”。这些材料可支持年序、官方决定、战事或特定地点的线索，却难以直接统计所有人的财富、迁徙、信仰和动机。更稳妥的读法是区分诏令与实施、军报与人口、行政称谓与自我认同，并让地方材料的缺环限制解释的范围。

由此可见，该事件不是孤立的年表点，而是资源、信息与权力关系重新配置的一环。不同地区的官员、社群和家庭可能以不同方式承担征发、保护、迁移或宗教活动；现存来源只能照亮其中一部分，不能替沉默者制造统一的经验。

\eventanchor{EH-0219-FANCHENG-GUAN-YU}{建安二十四年（219年）·关羽北攻樊城并水淹于禁军，后孙权袭取荆州，关羽兵败被杀}账本记录关羽集团、曹操集团与孙权集团在建安二十四年（219年）发生“关羽北攻樊城并水淹于禁军，后孙权袭取荆州，关羽兵败被杀”。条目把背景概括为关羽北进利用汉水洪灾压迫曹军，却使荆州后方防御薄弱；孙权长期不满荆州分配而转向突袭，并指出关羽集团覆灭，刘备失去荆州与东部通道，孙刘同盟实质破裂，曹魏、蜀汉、孙吴三方边界趋于定型。这个节点进入区域社会时，要经由文书、道路、军队、市场、寺院和家庭等不同中介；因此，政治行动的宣布、地方接收和日常生活的改变不能被视作同一时刻完成。

本条依据《三国志·关羽传》；《三国志·于禁传》；《三国志·吕蒙传》；《资治通鉴》卷六十八，证据提示为“战役序列可靠；水淹原因、吕蒙白衣渡江细节和关羽被杀地点有争议”。这些材料可支持年序、官方决定、战事或特定地点的线索，却难以直接统计所有人的财富、迁徙、信仰和动机。更稳妥的读法是区分诏令与实施、军报与人口、行政称谓与自我认同，并让地方材料的缺环限制解释的范围。

由此可见，该事件不是孤立的年表点，而是资源、信息与权力关系重新配置的一环。不同地区的官员、社群和家庭可能以不同方式承担征发、保护、迁移或宗教活动；现存来源只能照亮其中一部分，不能替沉默者制造统一的经验。

\eventanchor{EH-0220-CAO-CAO-DEATH}{建安二十五年正月（220年）·曹操去世，曹丕继承魏王位并控制汉廷军政机构}账本记录献帝朝廷与曹魏集团在建安二十五年正月（220年）发生“曹操去世，曹丕继承魏王位并控制汉廷军政机构”。条目把背景概括为曹操已以魏王国掌握北方资源和朝廷诏令，其死亡迫使曹氏迅速证明继承与军队控制的连续性，并指出曹丕顺利接管魏国，为迫使献帝禅位扫除最大内部障碍，东汉皇统仅余形式性地位。这个节点进入区域社会时，要经由文书、道路、军队、市场、寺院和家庭等不同中介；因此，政治行动的宣布、地方接收和日常生活的改变不能被视作同一时刻完成。

本条依据《三国志·武帝纪》；《三国志·文帝纪》；《后汉书·献帝纪》，证据提示为“卒年与继承事实可靠；遗令执行、继承协商和各地反应有争议”。这些材料可支持年序、官方决定、战事或特定地点的线索，却难以直接统计所有人的财富、迁徙、信仰和动机。更稳妥的读法是区分诏令与实施、军报与人口、行政称谓与自我认同，并让地方材料的缺环限制解释的范围。

由此可见，该事件不是孤立的年表点，而是资源、信息与权力关系重新配置的一环。不同地区的官员、社群和家庭可能以不同方式承担征发、保护、迁移或宗教活动；现存来源只能照亮其中一部分，不能替沉默者制造统一的经验。

\eventanchor{EH-0220-HAN-ABDICATION}{延康元年十月（220年）·献帝禅位于曹丕，曹丕建魏，东汉皇帝继承体系终结}账本记录东汉与曹魏在延康元年十月（220年）发生“献帝禅位于曹丕，曹丕建魏，东汉皇帝继承体系终结”。条目把背景概括为曹丕已继承魏王军政资源并控制献帝朝廷，以受禅礼仪将事实上的权力转移包装为合法改朝，并指出东汉作为中央王朝结束；蜀汉与孙吴不承认魏的统一权威，汉末危机转入三国并立与正统竞争。这个节点进入区域社会时，要经由文书、道路、军队、市场、寺院和家庭等不同中介；因此，政治行动的宣布、地方接收和日常生活的改变不能被视作同一时刻完成。

本条依据《后汉书·献帝纪》；《三国志·文帝纪》；《资治通鉴》卷六十九，证据提示为“禅代年月可靠；诏册撰作、群臣参与和“自愿”性质有争议”。这些材料可支持年序、官方决定、战事或特定地点的线索，却难以直接统计所有人的财富、迁徙、信仰和动机。更稳妥的读法是区分诏令与实施、军报与人口、行政称谓与自我认同，并让地方材料的缺环限制解释的范围。

由此可见，该事件不是孤立的年表点，而是资源、信息与权力关系重新配置的一环。不同地区的官员、社群和家庭可能以不同方式承担征发、保护、迁移或宗教活动；现存来源只能照亮其中一部分，不能替沉默者制造统一的经验。

\eventanchor{EH-LEDGER-001-1}{建武元年六月（25年）·“刘秀在鄗即皇帝位，建元建武，重建汉朝名号”的前期动员与资源准备}账本记录更始政权余部与刘秀集团在建武元年六月（25年）发生““刘秀在鄗即皇帝位，建元建武，重建汉朝名号”的前期动员与资源准备”。条目把背景概括为当时的权力、资源与信息约束推动相关过程展开，并指出这一过程为“刘秀在鄗即皇帝位，建元建武，重建汉朝名号”提供可观察的条件。这个节点进入区域社会时，要经由文书、道路、军队、市场、寺院和家庭等不同中介；因此，政治行动的宣布、地方接收和日常生活的改变不能被视作同一时刻完成。

本条依据《后汉书·光武帝纪上》；《后汉纪》卷一；《资治通鉴》卷四十，证据提示为“核心事项已有既有账本来源；本分解节点的参与者、执行范围和时间先后仍须逐案核验”。这些材料可支持年序、官方决定、战事或特定地点的线索，却难以直接统计所有人的财富、迁徙、信仰和动机。更稳妥的读法是区分诏令与实施、军报与人口、行政称谓与自我认同，并让地方材料的缺环限制解释的范围。

由此可见，该事件不是孤立的年表点，而是资源、信息与权力关系重新配置的一环。不同地区的官员、社群和家庭可能以不同方式承担征发、保护、迁移或宗教活动；现存来源只能照亮其中一部分，不能替沉默者制造统一的经验。

\eventanchor{EH-LEDGER-001-2}{建武元年六月（25年）·“刘秀在鄗即皇帝位，建元建武，重建汉朝名号”的命令、联盟或地方响应}账本记录更始政权余部与刘秀集团在建武元年六月（25年）发生““刘秀在鄗即皇帝位，建元建武，重建汉朝名号”的命令、联盟或地方响应”。条目把背景概括为当时的权力、资源与信息约束推动相关过程展开，并指出这一过程为“刘秀在鄗即皇帝位，建元建武，重建汉朝名号”提供可观察的条件。这个节点进入区域社会时，要经由文书、道路、军队、市场、寺院和家庭等不同中介；因此，政治行动的宣布、地方接收和日常生活的改变不能被视作同一时刻完成。

本条依据《后汉书·光武帝纪上》；《后汉纪》卷一；《资治通鉴》卷四十，证据提示为“核心事项已有既有账本来源；本分解节点的参与者、执行范围和时间先后仍须逐案核验”。这些材料可支持年序、官方决定、战事或特定地点的线索，却难以直接统计所有人的财富、迁徙、信仰和动机。更稳妥的读法是区分诏令与实施、军报与人口、行政称谓与自我认同，并让地方材料的缺环限制解释的范围。

由此可见，该事件不是孤立的年表点，而是资源、信息与权力关系重新配置的一环。不同地区的官员、社群和家庭可能以不同方式承担征发、保护、迁移或宗教活动；现存来源只能照亮其中一部分，不能替沉默者制造统一的经验。

\eventanchor{EH-LEDGER-001-3}{建武元年六月（25年）·刘秀在鄗即皇帝位，建元建武，重建汉朝名号}账本记录更始政权余部与刘秀集团在建武元年六月（25年）发生“刘秀在鄗即皇帝位，建元建武，重建汉朝名号”。条目把背景概括为“刘秀在鄗即皇帝位，建元建武，重建汉朝名号”发生的政治与区域条件共同作用，并指出该节点改变后续资源、联盟或制度处置的条件。这个节点进入区域社会时，要经由文书、道路、军队、市场、寺院和家庭等不同中介；因此，政治行动的宣布、地方接收和日常生活的改变不能被视作同一时刻完成。

本条依据《后汉书·光武帝纪上》；《后汉纪》卷一；《资治通鉴》卷四十，证据提示为“核心事项已有既有账本来源；本分解节点的参与者、执行范围和时间先后仍须逐案核验”。这些材料可支持年序、官方决定、战事或特定地点的线索，却难以直接统计所有人的财富、迁徙、信仰和动机。更稳妥的读法是区分诏令与实施、军报与人口、行政称谓与自我认同，并让地方材料的缺环限制解释的范围。

由此可见，该事件不是孤立的年表点，而是资源、信息与权力关系重新配置的一环。不同地区的官员、社群和家庭可能以不同方式承担征发、保护、迁移或宗教活动；现存来源只能照亮其中一部分，不能替沉默者制造统一的经验。

\eventanchor{EH-LEDGER-001-4}{建武元年六月（25年）·“刘秀在鄗即皇帝位，建元建武，重建汉朝名号”的执行中的空间与人员处置}账本记录更始政权余部与刘秀集团在建武元年六月（25年）发生““刘秀在鄗即皇帝位，建元建武，重建汉朝名号”的执行中的空间与人员处置”。条目把背景概括为“刘秀在鄗即皇帝位，建元建武，重建汉朝名号”发生的政治与区域条件共同作用，并指出该节点改变后续资源、联盟或制度处置的条件。这个节点进入区域社会时，要经由文书、道路、军队、市场、寺院和家庭等不同中介；因此，政治行动的宣布、地方接收和日常生活的改变不能被视作同一时刻完成。

本条依据《后汉书·光武帝纪上》；《后汉纪》卷一；《资治通鉴》卷四十，证据提示为“核心事项已有既有账本来源；本分解节点的参与者、执行范围和时间先后仍须逐案核验”。这些材料可支持年序、官方决定、战事或特定地点的线索，却难以直接统计所有人的财富、迁徙、信仰和动机。更稳妥的读法是区分诏令与实施、军报与人口、行政称谓与自我认同，并让地方材料的缺环限制解释的范围。

由此可见，该事件不是孤立的年表点，而是资源、信息与权力关系重新配置的一环。不同地区的官员、社群和家庭可能以不同方式承担征发、保护、迁移或宗教活动；现存来源只能照亮其中一部分，不能替沉默者制造统一的经验。

\eventanchor{EH-LEDGER-001-5}{建武元年六月（25年）·“刘秀在鄗即皇帝位，建元建武，重建汉朝名号”的后续制度或军政调整}账本记录更始政权余部与刘秀集团在建武元年六月（25年）发生““刘秀在鄗即皇帝位，建元建武，重建汉朝名号”的后续制度或军政调整”。条目把背景概括为核心节点发生后仍须处理其遗留问题，并指出后续影响须在不同地区和材料类型中分别核验。这个节点进入区域社会时，要经由文书、道路、军队、市场、寺院和家庭等不同中介；因此，政治行动的宣布、地方接收和日常生活的改变不能被视作同一时刻完成。

本条依据《后汉书·光武帝纪上》；《后汉纪》卷一；《资治通鉴》卷四十，证据提示为“核心事项已有既有账本来源；本分解节点的参与者、执行范围和时间先后仍须逐案核验”。这些材料可支持年序、官方决定、战事或特定地点的线索，却难以直接统计所有人的财富、迁徙、信仰和动机。更稳妥的读法是区分诏令与实施、军报与人口、行政称谓与自我认同，并让地方材料的缺环限制解释的范围。

由此可见，该事件不是孤立的年表点，而是资源、信息与权力关系重新配置的一环。不同地区的官员、社群和家庭可能以不同方式承担征发、保护、迁移或宗教活动；现存来源只能照亮其中一部分，不能替沉默者制造统一的经验。

\eventanchor{EH-LEDGER-001-6}{建武元年六月（25年）·“刘秀在鄗即皇帝位，建元建武，重建汉朝名号”的异源材料与影响范围的复核}账本记录更始政权余部与刘秀集团在建武元年六月（25年）发生““刘秀在鄗即皇帝位，建元建武，重建汉朝名号”的异源材料与影响范围的复核”。条目把背景概括为核心节点发生后仍须处理其遗留问题，并指出后续影响须在不同地区和材料类型中分别核验。这个节点进入区域社会时，要经由文书、道路、军队、市场、寺院和家庭等不同中介；因此，政治行动的宣布、地方接收和日常生活的改变不能被视作同一时刻完成。

本条依据《后汉书·光武帝纪上》；《后汉纪》卷一；《资治通鉴》卷四十，证据提示为“核心事项已有既有账本来源；本分解节点的参与者、执行范围和时间先后仍须逐案核验”。这些材料可支持年序、官方决定、战事或特定地点的线索，却难以直接统计所有人的财富、迁徙、信仰和动机。更稳妥的读法是区分诏令与实施、军报与人口、行政称谓与自我认同，并让地方材料的缺环限制解释的范围。

由此可见，该事件不是孤立的年表点，而是资源、信息与权力关系重新配置的一环。不同地区的官员、社群和家庭可能以不同方式承担征发、保护、迁移或宗教活动；现存来源只能照亮其中一部分，不能替沉默者制造统一的经验。

\eventanchor{EH-LEDGER-002-1}{建武元年（25年）·“刘秀定都洛阳，以雒阳为中央政府与军政调度中心”的前期动员与资源准备}账本记录东汉在建武元年（25年）发生““刘秀定都洛阳，以雒阳为中央政府与军政调度中心”的前期动员与资源准备”。条目把背景概括为当时的权力、资源与信息约束推动相关过程展开，并指出这一过程为“刘秀定都洛阳，以雒阳为中央政府与军政调度中心”提供可观察的条件。这个节点进入区域社会时，要经由文书、道路、军队、市场、寺院和家庭等不同中介；因此，政治行动的宣布、地方接收和日常生活的改变不能被视作同一时刻完成。

本条依据《后汉书·光武帝纪上》；《后汉纪》卷一；汉魏洛阳城遗址发掘报告，证据提示为“核心事项已有既有账本来源；本分解节点的参与者、执行范围和时间先后仍须逐案核验”。这些材料可支持年序、官方决定、战事或特定地点的线索，却难以直接统计所有人的财富、迁徙、信仰和动机。更稳妥的读法是区分诏令与实施、军报与人口、行政称谓与自我认同，并让地方材料的缺环限制解释的范围。

由此可见，该事件不是孤立的年表点，而是资源、信息与权力关系重新配置的一环。不同地区的官员、社群和家庭可能以不同方式承担征发、保护、迁移或宗教活动；现存来源只能照亮其中一部分，不能替沉默者制造统一的经验。

\eventanchor{EH-LEDGER-002-2}{建武元年（25年）·“刘秀定都洛阳，以雒阳为中央政府与军政调度中心”的命令、联盟或地方响应}账本记录东汉在建武元年（25年）发生““刘秀定都洛阳，以雒阳为中央政府与军政调度中心”的命令、联盟或地方响应”。条目把背景概括为当时的权力、资源与信息约束推动相关过程展开，并指出这一过程为“刘秀定都洛阳，以雒阳为中央政府与军政调度中心”提供可观察的条件。这个节点进入区域社会时，要经由文书、道路、军队、市场、寺院和家庭等不同中介；因此，政治行动的宣布、地方接收和日常生活的改变不能被视作同一时刻完成。

本条依据《后汉书·光武帝纪上》；《后汉纪》卷一；汉魏洛阳城遗址发掘报告，证据提示为“核心事项已有既有账本来源；本分解节点的参与者、执行范围和时间先后仍须逐案核验”。这些材料可支持年序、官方决定、战事或特定地点的线索，却难以直接统计所有人的财富、迁徙、信仰和动机。更稳妥的读法是区分诏令与实施、军报与人口、行政称谓与自我认同，并让地方材料的缺环限制解释的范围。

由此可见，该事件不是孤立的年表点，而是资源、信息与权力关系重新配置的一环。不同地区的官员、社群和家庭可能以不同方式承担征发、保护、迁移或宗教活动；现存来源只能照亮其中一部分，不能替沉默者制造统一的经验。

\eventanchor{EH-LEDGER-002-3}{建武元年（25年）·刘秀定都洛阳，以雒阳为中央政府与军政调度中心}账本记录东汉在建武元年（25年）发生“刘秀定都洛阳，以雒阳为中央政府与军政调度中心”。条目把背景概括为“刘秀定都洛阳，以雒阳为中央政府与军政调度中心”发生的政治与区域条件共同作用，并指出该节点改变后续资源、联盟或制度处置的条件。这个节点进入区域社会时，要经由文书、道路、军队、市场、寺院和家庭等不同中介；因此，政治行动的宣布、地方接收和日常生活的改变不能被视作同一时刻完成。

本条依据《后汉书·光武帝纪上》；《后汉纪》卷一；汉魏洛阳城遗址发掘报告，证据提示为“核心事项已有既有账本来源；本分解节点的参与者、执行范围和时间先后仍须逐案核验”。这些材料可支持年序、官方决定、战事或特定地点的线索，却难以直接统计所有人的财富、迁徙、信仰和动机。更稳妥的读法是区分诏令与实施、军报与人口、行政称谓与自我认同，并让地方材料的缺环限制解释的范围。

由此可见，该事件不是孤立的年表点，而是资源、信息与权力关系重新配置的一环。不同地区的官员、社群和家庭可能以不同方式承担征发、保护、迁移或宗教活动；现存来源只能照亮其中一部分，不能替沉默者制造统一的经验。

\eventanchor{EH-LEDGER-002-4}{建武元年（25年）·“刘秀定都洛阳，以雒阳为中央政府与军政调度中心”的执行中的空间与人员处置}账本记录东汉在建武元年（25年）发生““刘秀定都洛阳，以雒阳为中央政府与军政调度中心”的执行中的空间与人员处置”。条目把背景概括为“刘秀定都洛阳，以雒阳为中央政府与军政调度中心”发生的政治与区域条件共同作用，并指出该节点改变后续资源、联盟或制度处置的条件。这个节点进入区域社会时，要经由文书、道路、军队、市场、寺院和家庭等不同中介；因此，政治行动的宣布、地方接收和日常生活的改变不能被视作同一时刻完成。

本条依据《后汉书·光武帝纪上》；《后汉纪》卷一；汉魏洛阳城遗址发掘报告，证据提示为“核心事项已有既有账本来源；本分解节点的参与者、执行范围和时间先后仍须逐案核验”。这些材料可支持年序、官方决定、战事或特定地点的线索，却难以直接统计所有人的财富、迁徙、信仰和动机。更稳妥的读法是区分诏令与实施、军报与人口、行政称谓与自我认同，并让地方材料的缺环限制解释的范围。

由此可见，该事件不是孤立的年表点，而是资源、信息与权力关系重新配置的一环。不同地区的官员、社群和家庭可能以不同方式承担征发、保护、迁移或宗教活动；现存来源只能照亮其中一部分，不能替沉默者制造统一的经验。

\eventanchor{EH-LEDGER-002-5}{建武元年（25年）·“刘秀定都洛阳，以雒阳为中央政府与军政调度中心”的后续制度或军政调整}账本记录东汉在建武元年（25年）发生““刘秀定都洛阳，以雒阳为中央政府与军政调度中心”的后续制度或军政调整”。条目把背景概括为核心节点发生后仍须处理其遗留问题，并指出后续影响须在不同地区和材料类型中分别核验。这个节点进入区域社会时，要经由文书、道路、军队、市场、寺院和家庭等不同中介；因此，政治行动的宣布、地方接收和日常生活的改变不能被视作同一时刻完成。

本条依据《后汉书·光武帝纪上》；《后汉纪》卷一；汉魏洛阳城遗址发掘报告，证据提示为“核心事项已有既有账本来源；本分解节点的参与者、执行范围和时间先后仍须逐案核验”。这些材料可支持年序、官方决定、战事或特定地点的线索，却难以直接统计所有人的财富、迁徙、信仰和动机。更稳妥的读法是区分诏令与实施、军报与人口、行政称谓与自我认同，并让地方材料的缺环限制解释的范围。

由此可见，该事件不是孤立的年表点，而是资源、信息与权力关系重新配置的一环。不同地区的官员、社群和家庭可能以不同方式承担征发、保护、迁移或宗教活动；现存来源只能照亮其中一部分，不能替沉默者制造统一的经验。

\eventanchor{EH-LEDGER-002-6}{建武元年（25年）·“刘秀定都洛阳，以雒阳为中央政府与军政调度中心”的异源材料与影响范围的复核}账本记录东汉在建武元年（25年）发生““刘秀定都洛阳，以雒阳为中央政府与军政调度中心”的异源材料与影响范围的复核”。条目把背景概括为核心节点发生后仍须处理其遗留问题，并指出后续影响须在不同地区和材料类型中分别核验。这个节点进入区域社会时，要经由文书、道路、军队、市场、寺院和家庭等不同中介；因此，政治行动的宣布、地方接收和日常生活的改变不能被视作同一时刻完成。

本条依据《后汉书·光武帝纪上》；《后汉纪》卷一；汉魏洛阳城遗址发掘报告，证据提示为“核心事项已有既有账本来源；本分解节点的参与者、执行范围和时间先后仍须逐案核验”。这些材料可支持年序、官方决定、战事或特定地点的线索，却难以直接统计所有人的财富、迁徙、信仰和动机。更稳妥的读法是区分诏令与实施、军报与人口、行政称谓与自我认同，并让地方材料的缺环限制解释的范围。

由此可见，该事件不是孤立的年表点，而是资源、信息与权力关系重新配置的一环。不同地区的官员、社群和家庭可能以不同方式承担征发、保护、迁移或宗教活动；现存来源只能照亮其中一部分，不能替沉默者制造统一的经验。

\eventanchor{EH-LEDGER-003-1}{建武二年（26年）·“赤眉军入长安，更始帝刘玄败亡，关中更始政权终结”的前期动员与资源准备}账本记录赤眉军与更始政权在建武二年（26年）发生““赤眉军入长安，更始帝刘玄败亡，关中更始政权终结”的前期动员与资源准备”。条目把背景概括为当时的权力、资源与信息约束推动相关过程展开，并指出这一过程为“赤眉军入长安，更始帝刘玄败亡，关中更始政权终结”提供可观察的条件。这个节点进入区域社会时，要经由文书、道路、军队、市场、寺院和家庭等不同中介；因此，政治行动的宣布、地方接收和日常生活的改变不能被视作同一时刻完成。

本条依据《后汉书·刘玄传》；《后汉纪》卷二；《资治通鉴》卷四十，证据提示为“核心事项已有既有账本来源；本分解节点的参与者、执行范围和时间先后仍须逐案核验”。这些材料可支持年序、官方决定、战事或特定地点的线索，却难以直接统计所有人的财富、迁徙、信仰和动机。更稳妥的读法是区分诏令与实施、军报与人口、行政称谓与自我认同，并让地方材料的缺环限制解释的范围。

由此可见，该事件不是孤立的年表点，而是资源、信息与权力关系重新配置的一环。不同地区的官员、社群和家庭可能以不同方式承担征发、保护、迁移或宗教活动；现存来源只能照亮其中一部分，不能替沉默者制造统一的经验。

\subsection{社会关系与宗教空间}

\eventanchor{EH-LEDGER-003-2}{建武二年（26年）·“赤眉军入长安，更始帝刘玄败亡，关中更始政权终结”的命令、联盟或地方响应}账本记录赤眉军与更始政权在建武二年（26年）发生““赤眉军入长安，更始帝刘玄败亡，关中更始政权终结”的命令、联盟或地方响应”。条目把背景概括为当时的权力、资源与信息约束推动相关过程展开，并指出这一过程为“赤眉军入长安，更始帝刘玄败亡，关中更始政权终结”提供可观察的条件。这个节点进入区域社会时，要经由文书、道路、军队、市场、寺院和家庭等不同中介；因此，政治行动的宣布、地方接收和日常生活的改变不能被视作同一时刻完成。

本条依据《后汉书·刘玄传》；《后汉纪》卷二；《资治通鉴》卷四十，证据提示为“核心事项已有既有账本来源；本分解节点的参与者、执行范围和时间先后仍须逐案核验”。这些材料可支持年序、官方决定、战事或特定地点的线索，却难以直接统计所有人的财富、迁徙、信仰和动机。更稳妥的读法是区分诏令与实施、军报与人口、行政称谓与自我认同，并让地方材料的缺环限制解释的范围。

由此可见，该事件不是孤立的年表点，而是资源、信息与权力关系重新配置的一环。不同地区的官员、社群和家庭可能以不同方式承担征发、保护、迁移或宗教活动；现存来源只能照亮其中一部分，不能替沉默者制造统一的经验。

\eventanchor{EH-LEDGER-003-3}{建武二年（26年）·赤眉军入长安，更始帝刘玄败亡，关中更始政权终结}账本记录赤眉军与更始政权在建武二年（26年）发生“赤眉军入长安，更始帝刘玄败亡，关中更始政权终结”。条目把背景概括为“赤眉军入长安，更始帝刘玄败亡，关中更始政权终结”发生的政治与区域条件共同作用，并指出该节点改变后续资源、联盟或制度处置的条件。这个节点进入区域社会时，要经由文书、道路、军队、市场、寺院和家庭等不同中介；因此，政治行动的宣布、地方接收和日常生活的改变不能被视作同一时刻完成。

本条依据《后汉书·刘玄传》；《后汉纪》卷二；《资治通鉴》卷四十，证据提示为“核心事项已有既有账本来源；本分解节点的参与者、执行范围和时间先后仍须逐案核验”。这些材料可支持年序、官方决定、战事或特定地点的线索，却难以直接统计所有人的财富、迁徙、信仰和动机。更稳妥的读法是区分诏令与实施、军报与人口、行政称谓与自我认同，并让地方材料的缺环限制解释的范围。

由此可见，该事件不是孤立的年表点，而是资源、信息与权力关系重新配置的一环。不同地区的官员、社群和家庭可能以不同方式承担征发、保护、迁移或宗教活动；现存来源只能照亮其中一部分，不能替沉默者制造统一的经验。

\eventanchor{EH-LEDGER-003-4}{建武二年（26年）·“赤眉军入长安，更始帝刘玄败亡，关中更始政权终结”的执行中的空间与人员处置}账本记录赤眉军与更始政权在建武二年（26年）发生““赤眉军入长安，更始帝刘玄败亡，关中更始政权终结”的执行中的空间与人员处置”。条目把背景概括为“赤眉军入长安，更始帝刘玄败亡，关中更始政权终结”发生的政治与区域条件共同作用，并指出该节点改变后续资源、联盟或制度处置的条件。这个节点进入区域社会时，要经由文书、道路、军队、市场、寺院和家庭等不同中介；因此，政治行动的宣布、地方接收和日常生活的改变不能被视作同一时刻完成。

本条依据《后汉书·刘玄传》；《后汉纪》卷二；《资治通鉴》卷四十，证据提示为“核心事项已有既有账本来源；本分解节点的参与者、执行范围和时间先后仍须逐案核验”。这些材料可支持年序、官方决定、战事或特定地点的线索，却难以直接统计所有人的财富、迁徙、信仰和动机。更稳妥的读法是区分诏令与实施、军报与人口、行政称谓与自我认同，并让地方材料的缺环限制解释的范围。

由此可见，该事件不是孤立的年表点，而是资源、信息与权力关系重新配置的一环。不同地区的官员、社群和家庭可能以不同方式承担征发、保护、迁移或宗教活动；现存来源只能照亮其中一部分，不能替沉默者制造统一的经验。

\eventanchor{EH-LEDGER-003-5}{建武二年（26年）·“赤眉军入长安，更始帝刘玄败亡，关中更始政权终结”的后续制度或军政调整}账本记录赤眉军与更始政权在建武二年（26年）发生““赤眉军入长安，更始帝刘玄败亡，关中更始政权终结”的后续制度或军政调整”。条目把背景概括为核心节点发生后仍须处理其遗留问题，并指出后续影响须在不同地区和材料类型中分别核验。这个节点进入区域社会时，要经由文书、道路、军队、市场、寺院和家庭等不同中介；因此，政治行动的宣布、地方接收和日常生活的改变不能被视作同一时刻完成。

本条依据《后汉书·刘玄传》；《后汉纪》卷二；《资治通鉴》卷四十，证据提示为“核心事项已有既有账本来源；本分解节点的参与者、执行范围和时间先后仍须逐案核验”。这些材料可支持年序、官方决定、战事或特定地点的线索，却难以直接统计所有人的财富、迁徙、信仰和动机。更稳妥的读法是区分诏令与实施、军报与人口、行政称谓与自我认同，并让地方材料的缺环限制解释的范围。

由此可见，该事件不是孤立的年表点，而是资源、信息与权力关系重新配置的一环。不同地区的官员、社群和家庭可能以不同方式承担征发、保护、迁移或宗教活动；现存来源只能照亮其中一部分，不能替沉默者制造统一的经验。

\eventanchor{EH-LEDGER-003-6}{建武二年（26年）·“赤眉军入长安，更始帝刘玄败亡，关中更始政权终结”的异源材料与影响范围的复核}账本记录赤眉军与更始政权在建武二年（26年）发生““赤眉军入长安，更始帝刘玄败亡，关中更始政权终结”的异源材料与影响范围的复核”。条目把背景概括为核心节点发生后仍须处理其遗留问题，并指出后续影响须在不同地区和材料类型中分别核验。这个节点进入区域社会时，要经由文书、道路、军队、市场、寺院和家庭等不同中介；因此，政治行动的宣布、地方接收和日常生活的改变不能被视作同一时刻完成。

本条依据《后汉书·刘玄传》；《后汉纪》卷二；《资治通鉴》卷四十，证据提示为“核心事项已有既有账本来源；本分解节点的参与者、执行范围和时间先后仍须逐案核验”。这些材料可支持年序、官方决定、战事或特定地点的线索，却难以直接统计所有人的财富、迁徙、信仰和动机。更稳妥的读法是区分诏令与实施、军报与人口、行政称谓与自我认同，并让地方材料的缺环限制解释的范围。

由此可见，该事件不是孤立的年表点，而是资源、信息与权力关系重新配置的一环。不同地区的官员、社群和家庭可能以不同方式承担征发、保护、迁移或宗教活动；现存来源只能照亮其中一部分，不能替沉默者制造统一的经验。

\eventanchor{EH-LEDGER-004-1}{建武三年（27年）·“刘秀军在崤底一带受降赤眉，赤眉主力瓦解”的前期动员与资源准备}账本记录东汉与赤眉军在建武三年（27年）发生““刘秀军在崤底一带受降赤眉，赤眉主力瓦解”的前期动员与资源准备”。条目把背景概括为当时的权力、资源与信息约束推动相关过程展开，并指出这一过程为“刘秀军在崤底一带受降赤眉，赤眉主力瓦解”提供可观察的条件。这个节点进入区域社会时，要经由文书、道路、军队、市场、寺院和家庭等不同中介；因此，政治行动的宣布、地方接收和日常生活的改变不能被视作同一时刻完成。

本条依据《后汉书·光武帝纪上》；《后汉书·刘盆子传》；《后汉纪》卷二，证据提示为“核心事项已有既有账本来源；本分解节点的参与者、执行范围和时间先后仍须逐案核验”。这些材料可支持年序、官方决定、战事或特定地点的线索，却难以直接统计所有人的财富、迁徙、信仰和动机。更稳妥的读法是区分诏令与实施、军报与人口、行政称谓与自我认同，并让地方材料的缺环限制解释的范围。

由此可见，该事件不是孤立的年表点，而是资源、信息与权力关系重新配置的一环。不同地区的官员、社群和家庭可能以不同方式承担征发、保护、迁移或宗教活动；现存来源只能照亮其中一部分，不能替沉默者制造统一的经验。

\eventanchor{EH-LEDGER-004-2}{建武三年（27年）·“刘秀军在崤底一带受降赤眉，赤眉主力瓦解”的命令、联盟或地方响应}账本记录东汉与赤眉军在建武三年（27年）发生““刘秀军在崤底一带受降赤眉，赤眉主力瓦解”的命令、联盟或地方响应”。条目把背景概括为当时的权力、资源与信息约束推动相关过程展开，并指出这一过程为“刘秀军在崤底一带受降赤眉，赤眉主力瓦解”提供可观察的条件。这个节点进入区域社会时，要经由文书、道路、军队、市场、寺院和家庭等不同中介；因此，政治行动的宣布、地方接收和日常生活的改变不能被视作同一时刻完成。

本条依据《后汉书·光武帝纪上》；《后汉书·刘盆子传》；《后汉纪》卷二，证据提示为“核心事项已有既有账本来源；本分解节点的参与者、执行范围和时间先后仍须逐案核验”。这些材料可支持年序、官方决定、战事或特定地点的线索，却难以直接统计所有人的财富、迁徙、信仰和动机。更稳妥的读法是区分诏令与实施、军报与人口、行政称谓与自我认同，并让地方材料的缺环限制解释的范围。

由此可见，该事件不是孤立的年表点，而是资源、信息与权力关系重新配置的一环。不同地区的官员、社群和家庭可能以不同方式承担征发、保护、迁移或宗教活动；现存来源只能照亮其中一部分，不能替沉默者制造统一的经验。

\eventanchor{EH-LEDGER-004-3}{建武三年（27年）·刘秀军在崤底一带受降赤眉，赤眉主力瓦解}账本记录东汉与赤眉军在建武三年（27年）发生“刘秀军在崤底一带受降赤眉，赤眉主力瓦解”。条目把背景概括为“刘秀军在崤底一带受降赤眉，赤眉主力瓦解”发生的政治与区域条件共同作用，并指出该节点改变后续资源、联盟或制度处置的条件。这个节点进入区域社会时，要经由文书、道路、军队、市场、寺院和家庭等不同中介；因此，政治行动的宣布、地方接收和日常生活的改变不能被视作同一时刻完成。

本条依据《后汉书·光武帝纪上》；《后汉书·刘盆子传》；《后汉纪》卷二，证据提示为“核心事项已有既有账本来源；本分解节点的参与者、执行范围和时间先后仍须逐案核验”。这些材料可支持年序、官方决定、战事或特定地点的线索，却难以直接统计所有人的财富、迁徙、信仰和动机。更稳妥的读法是区分诏令与实施、军报与人口、行政称谓与自我认同，并让地方材料的缺环限制解释的范围。

由此可见，该事件不是孤立的年表点，而是资源、信息与权力关系重新配置的一环。不同地区的官员、社群和家庭可能以不同方式承担征发、保护、迁移或宗教活动；现存来源只能照亮其中一部分，不能替沉默者制造统一的经验。

\eventanchor{EH-LEDGER-004-4}{建武三年（27年）·“刘秀军在崤底一带受降赤眉，赤眉主力瓦解”的执行中的空间与人员处置}账本记录东汉与赤眉军在建武三年（27年）发生““刘秀军在崤底一带受降赤眉，赤眉主力瓦解”的执行中的空间与人员处置”。条目把背景概括为“刘秀军在崤底一带受降赤眉，赤眉主力瓦解”发生的政治与区域条件共同作用，并指出该节点改变后续资源、联盟或制度处置的条件。这个节点进入区域社会时，要经由文书、道路、军队、市场、寺院和家庭等不同中介；因此，政治行动的宣布、地方接收和日常生活的改变不能被视作同一时刻完成。

本条依据《后汉书·光武帝纪上》；《后汉书·刘盆子传》；《后汉纪》卷二，证据提示为“核心事项已有既有账本来源；本分解节点的参与者、执行范围和时间先后仍须逐案核验”。这些材料可支持年序、官方决定、战事或特定地点的线索，却难以直接统计所有人的财富、迁徙、信仰和动机。更稳妥的读法是区分诏令与实施、军报与人口、行政称谓与自我认同，并让地方材料的缺环限制解释的范围。

由此可见，该事件不是孤立的年表点，而是资源、信息与权力关系重新配置的一环。不同地区的官员、社群和家庭可能以不同方式承担征发、保护、迁移或宗教活动；现存来源只能照亮其中一部分，不能替沉默者制造统一的经验。

\eventanchor{EH-LEDGER-004-5}{建武三年（27年）·“刘秀军在崤底一带受降赤眉，赤眉主力瓦解”的后续制度或军政调整}账本记录东汉与赤眉军在建武三年（27年）发生““刘秀军在崤底一带受降赤眉，赤眉主力瓦解”的后续制度或军政调整”。条目把背景概括为核心节点发生后仍须处理其遗留问题，并指出后续影响须在不同地区和材料类型中分别核验。这个节点进入区域社会时，要经由文书、道路、军队、市场、寺院和家庭等不同中介；因此，政治行动的宣布、地方接收和日常生活的改变不能被视作同一时刻完成。

本条依据《后汉书·光武帝纪上》；《后汉书·刘盆子传》；《后汉纪》卷二，证据提示为“核心事项已有既有账本来源；本分解节点的参与者、执行范围和时间先后仍须逐案核验”。这些材料可支持年序、官方决定、战事或特定地点的线索，却难以直接统计所有人的财富、迁徙、信仰和动机。更稳妥的读法是区分诏令与实施、军报与人口、行政称谓与自我认同，并让地方材料的缺环限制解释的范围。

由此可见，该事件不是孤立的年表点，而是资源、信息与权力关系重新配置的一环。不同地区的官员、社群和家庭可能以不同方式承担征发、保护、迁移或宗教活动；现存来源只能照亮其中一部分，不能替沉默者制造统一的经验。

\eventanchor{EH-LEDGER-004-6}{建武三年（27年）·“刘秀军在崤底一带受降赤眉，赤眉主力瓦解”的异源材料与影响范围的复核}账本记录东汉与赤眉军在建武三年（27年）发生““刘秀军在崤底一带受降赤眉，赤眉主力瓦解”的异源材料与影响范围的复核”。条目把背景概括为核心节点发生后仍须处理其遗留问题，并指出后续影响须在不同地区和材料类型中分别核验。这个节点进入区域社会时，要经由文书、道路、军队、市场、寺院和家庭等不同中介；因此，政治行动的宣布、地方接收和日常生活的改变不能被视作同一时刻完成。

本条依据《后汉书·光武帝纪上》；《后汉书·刘盆子传》；《后汉纪》卷二，证据提示为“核心事项已有既有账本来源；本分解节点的参与者、执行范围和时间先后仍须逐案核验”。这些材料可支持年序、官方决定、战事或特定地点的线索，却难以直接统计所有人的财富、迁徙、信仰和动机。更稳妥的读法是区分诏令与实施、军报与人口、行政称谓与自我认同，并让地方材料的缺环限制解释的范围。

由此可见，该事件不是孤立的年表点，而是资源、信息与权力关系重新配置的一环。不同地区的官员、社群和家庭可能以不同方式承担征发、保护、迁移或宗教活动；现存来源只能照亮其中一部分，不能替沉默者制造统一的经验。

\eventanchor{EH-LEDGER-005-1}{建武六年（30年）·“公孙述在成都称帝，国号成家，巴蜀形成独立政权中心”的前期动员与资源准备}账本记录成家政权在建武六年（30年）发生““公孙述在成都称帝，国号成家，巴蜀形成独立政权中心”的前期动员与资源准备”。条目把背景概括为当时的权力、资源与信息约束推动相关过程展开，并指出这一过程为“公孙述在成都称帝，国号成家，巴蜀形成独立政权中心”提供可观察的条件。这个节点进入区域社会时，要经由文书、道路、军队、市场、寺院和家庭等不同中介；因此，政治行动的宣布、地方接收和日常生活的改变不能被视作同一时刻完成。

本条依据《后汉书·公孙述传》；《华阳国志·公孙述志》；《后汉纪》卷五，证据提示为“核心事项已有既有账本来源；本分解节点的参与者、执行范围和时间先后仍须逐案核验”。这些材料可支持年序、官方决定、战事或特定地点的线索，却难以直接统计所有人的财富、迁徙、信仰和动机。更稳妥的读法是区分诏令与实施、军报与人口、行政称谓与自我认同，并让地方材料的缺环限制解释的范围。

由此可见，该事件不是孤立的年表点，而是资源、信息与权力关系重新配置的一环。不同地区的官员、社群和家庭可能以不同方式承担征发、保护、迁移或宗教活动；现存来源只能照亮其中一部分，不能替沉默者制造统一的经验。

\eventanchor{EH-LEDGER-005-2}{建武六年（30年）·“公孙述在成都称帝，国号成家，巴蜀形成独立政权中心”的命令、联盟或地方响应}账本记录成家政权在建武六年（30年）发生““公孙述在成都称帝，国号成家，巴蜀形成独立政权中心”的命令、联盟或地方响应”。条目把背景概括为当时的权力、资源与信息约束推动相关过程展开，并指出这一过程为“公孙述在成都称帝，国号成家，巴蜀形成独立政权中心”提供可观察的条件。这个节点进入区域社会时，要经由文书、道路、军队、市场、寺院和家庭等不同中介；因此，政治行动的宣布、地方接收和日常生活的改变不能被视作同一时刻完成。

本条依据《后汉书·公孙述传》；《华阳国志·公孙述志》；《后汉纪》卷五，证据提示为“核心事项已有既有账本来源；本分解节点的参与者、执行范围和时间先后仍须逐案核验”。这些材料可支持年序、官方决定、战事或特定地点的线索，却难以直接统计所有人的财富、迁徙、信仰和动机。更稳妥的读法是区分诏令与实施、军报与人口、行政称谓与自我认同，并让地方材料的缺环限制解释的范围。

由此可见，该事件不是孤立的年表点，而是资源、信息与权力关系重新配置的一环。不同地区的官员、社群和家庭可能以不同方式承担征发、保护、迁移或宗教活动；现存来源只能照亮其中一部分，不能替沉默者制造统一的经验。

\eventanchor{EH-LEDGER-005-3}{建武六年（30年）·公孙述在成都称帝，国号成家，巴蜀形成独立政权中心}账本记录成家政权在建武六年（30年）发生“公孙述在成都称帝，国号成家，巴蜀形成独立政权中心”。条目把背景概括为“公孙述在成都称帝，国号成家，巴蜀形成独立政权中心”发生的政治与区域条件共同作用，并指出该节点改变后续资源、联盟或制度处置的条件。这个节点进入区域社会时，要经由文书、道路、军队、市场、寺院和家庭等不同中介；因此，政治行动的宣布、地方接收和日常生活的改变不能被视作同一时刻完成。

本条依据《后汉书·公孙述传》；《华阳国志·公孙述志》；《后汉纪》卷五，证据提示为“核心事项已有既有账本来源；本分解节点的参与者、执行范围和时间先后仍须逐案核验”。这些材料可支持年序、官方决定、战事或特定地点的线索，却难以直接统计所有人的财富、迁徙、信仰和动机。更稳妥的读法是区分诏令与实施、军报与人口、行政称谓与自我认同，并让地方材料的缺环限制解释的范围。

由此可见，该事件不是孤立的年表点，而是资源、信息与权力关系重新配置的一环。不同地区的官员、社群和家庭可能以不同方式承担征发、保护、迁移或宗教活动；现存来源只能照亮其中一部分，不能替沉默者制造统一的经验。

\eventanchor{EH-LEDGER-005-4}{建武六年（30年）·“公孙述在成都称帝，国号成家，巴蜀形成独立政权中心”的执行中的空间与人员处置}账本记录成家政权在建武六年（30年）发生““公孙述在成都称帝，国号成家，巴蜀形成独立政权中心”的执行中的空间与人员处置”。条目把背景概括为“公孙述在成都称帝，国号成家，巴蜀形成独立政权中心”发生的政治与区域条件共同作用，并指出该节点改变后续资源、联盟或制度处置的条件。这个节点进入区域社会时，要经由文书、道路、军队、市场、寺院和家庭等不同中介；因此，政治行动的宣布、地方接收和日常生活的改变不能被视作同一时刻完成。

本条依据《后汉书·公孙述传》；《华阳国志·公孙述志》；《后汉纪》卷五，证据提示为“核心事项已有既有账本来源；本分解节点的参与者、执行范围和时间先后仍须逐案核验”。这些材料可支持年序、官方决定、战事或特定地点的线索，却难以直接统计所有人的财富、迁徙、信仰和动机。更稳妥的读法是区分诏令与实施、军报与人口、行政称谓与自我认同，并让地方材料的缺环限制解释的范围。

由此可见，该事件不是孤立的年表点，而是资源、信息与权力关系重新配置的一环。不同地区的官员、社群和家庭可能以不同方式承担征发、保护、迁移或宗教活动；现存来源只能照亮其中一部分，不能替沉默者制造统一的经验。

\eventanchor{EH-LEDGER-005-5}{建武六年（30年）·“公孙述在成都称帝，国号成家，巴蜀形成独立政权中心”的后续制度或军政调整}账本记录成家政权在建武六年（30年）发生““公孙述在成都称帝，国号成家，巴蜀形成独立政权中心”的后续制度或军政调整”。条目把背景概括为核心节点发生后仍须处理其遗留问题，并指出后续影响须在不同地区和材料类型中分别核验。这个节点进入区域社会时，要经由文书、道路、军队、市场、寺院和家庭等不同中介；因此，政治行动的宣布、地方接收和日常生活的改变不能被视作同一时刻完成。

本条依据《后汉书·公孙述传》；《华阳国志·公孙述志》；《后汉纪》卷五，证据提示为“核心事项已有既有账本来源；本分解节点的参与者、执行范围和时间先后仍须逐案核验”。这些材料可支持年序、官方决定、战事或特定地点的线索，却难以直接统计所有人的财富、迁徙、信仰和动机。更稳妥的读法是区分诏令与实施、军报与人口、行政称谓与自我认同，并让地方材料的缺环限制解释的范围。

由此可见，该事件不是孤立的年表点，而是资源、信息与权力关系重新配置的一环。不同地区的官员、社群和家庭可能以不同方式承担征发、保护、迁移或宗教活动；现存来源只能照亮其中一部分，不能替沉默者制造统一的经验。

\eventanchor{EH-LEDGER-005-6}{建武六年（30年）·“公孙述在成都称帝，国号成家，巴蜀形成独立政权中心”的异源材料与影响范围的复核}账本记录成家政权在建武六年（30年）发生““公孙述在成都称帝，国号成家，巴蜀形成独立政权中心”的异源材料与影响范围的复核”。条目把背景概括为核心节点发生后仍须处理其遗留问题，并指出后续影响须在不同地区和材料类型中分别核验。这个节点进入区域社会时，要经由文书、道路、军队、市场、寺院和家庭等不同中介；因此，政治行动的宣布、地方接收和日常生活的改变不能被视作同一时刻完成。

本条依据《后汉书·公孙述传》；《华阳国志·公孙述志》；《后汉纪》卷五，证据提示为“核心事项已有既有账本来源；本分解节点的参与者、执行范围和时间先后仍须逐案核验”。这些材料可支持年序、官方决定、战事或特定地点的线索，却难以直接统计所有人的财富、迁徙、信仰和动机。更稳妥的读法是区分诏令与实施、军报与人口、行政称谓与自我认同，并让地方材料的缺环限制解释的范围。

由此可见，该事件不是孤立的年表点，而是资源、信息与权力关系重新配置的一环。不同地区的官员、社群和家庭可能以不同方式承担征发、保护、迁移或宗教活动；现存来源只能照亮其中一部分，不能替沉默者制造统一的经验。

\eventanchor{EH-LEDGER-006-1}{建武十年（34年）·“隗嚣死后，其子隗纯降汉，陇右主要割据力量被解除”的前期动员与资源准备}账本记录东汉与陇右集团在建武十年（34年）发生““隗嚣死后，其子隗纯降汉，陇右主要割据力量被解除”的前期动员与资源准备”。条目把背景概括为当时的权力、资源与信息约束推动相关过程展开，并指出这一过程为“隗嚣死后，其子隗纯降汉，陇右主要割据力量被解除”提供可观察的条件。这个节点进入区域社会时，要经由文书、道路、军队、市场、寺院和家庭等不同中介；因此，政治行动的宣布、地方接收和日常生活的改变不能被视作同一时刻完成。

本条依据《后汉书·隗嚣传》；《后汉纪》卷八；《资治通鉴》卷四十三，证据提示为“核心事项已有既有账本来源；本分解节点的参与者、执行范围和时间先后仍须逐案核验”。这些材料可支持年序、官方决定、战事或特定地点的线索，却难以直接统计所有人的财富、迁徙、信仰和动机。更稳妥的读法是区分诏令与实施、军报与人口、行政称谓与自我认同，并让地方材料的缺环限制解释的范围。

由此可见，该事件不是孤立的年表点，而是资源、信息与权力关系重新配置的一环。不同地区的官员、社群和家庭可能以不同方式承担征发、保护、迁移或宗教活动；现存来源只能照亮其中一部分，不能替沉默者制造统一的经验。

\eventanchor{EH-LEDGER-006-2}{建武十年（34年）·“隗嚣死后，其子隗纯降汉，陇右主要割据力量被解除”的命令、联盟或地方响应}账本记录东汉与陇右集团在建武十年（34年）发生““隗嚣死后，其子隗纯降汉，陇右主要割据力量被解除”的命令、联盟或地方响应”。条目把背景概括为当时的权力、资源与信息约束推动相关过程展开，并指出这一过程为“隗嚣死后，其子隗纯降汉，陇右主要割据力量被解除”提供可观察的条件。这个节点进入区域社会时，要经由文书、道路、军队、市场、寺院和家庭等不同中介；因此，政治行动的宣布、地方接收和日常生活的改变不能被视作同一时刻完成。

本条依据《后汉书·隗嚣传》；《后汉纪》卷八；《资治通鉴》卷四十三，证据提示为“核心事项已有既有账本来源；本分解节点的参与者、执行范围和时间先后仍须逐案核验”。这些材料可支持年序、官方决定、战事或特定地点的线索，却难以直接统计所有人的财富、迁徙、信仰和动机。更稳妥的读法是区分诏令与实施、军报与人口、行政称谓与自我认同，并让地方材料的缺环限制解释的范围。

由此可见，该事件不是孤立的年表点，而是资源、信息与权力关系重新配置的一环。不同地区的官员、社群和家庭可能以不同方式承担征发、保护、迁移或宗教活动；现存来源只能照亮其中一部分，不能替沉默者制造统一的经验。

\eventanchor{EH-LEDGER-006-3}{建武十年（34年）·隗嚣死后，其子隗纯降汉，陇右主要割据力量被解除}账本记录东汉与陇右集团在建武十年（34年）发生“隗嚣死后，其子隗纯降汉，陇右主要割据力量被解除”。条目把背景概括为“隗嚣死后，其子隗纯降汉，陇右主要割据力量被解除”发生的政治与区域条件共同作用，并指出该节点改变后续资源、联盟或制度处置的条件。这个节点进入区域社会时，要经由文书、道路、军队、市场、寺院和家庭等不同中介；因此，政治行动的宣布、地方接收和日常生活的改变不能被视作同一时刻完成。

本条依据《后汉书·隗嚣传》；《后汉纪》卷八；《资治通鉴》卷四十三，证据提示为“核心事项已有既有账本来源；本分解节点的参与者、执行范围和时间先后仍须逐案核验”。这些材料可支持年序、官方决定、战事或特定地点的线索，却难以直接统计所有人的财富、迁徙、信仰和动机。更稳妥的读法是区分诏令与实施、军报与人口、行政称谓与自我认同，并让地方材料的缺环限制解释的范围。

由此可见，该事件不是孤立的年表点，而是资源、信息与权力关系重新配置的一环。不同地区的官员、社群和家庭可能以不同方式承担征发、保护、迁移或宗教活动；现存来源只能照亮其中一部分，不能替沉默者制造统一的经验。

\eventanchor{EH-LEDGER-006-4}{建武十年（34年）·“隗嚣死后，其子隗纯降汉，陇右主要割据力量被解除”的执行中的空间与人员处置}账本记录东汉与陇右集团在建武十年（34年）发生““隗嚣死后，其子隗纯降汉，陇右主要割据力量被解除”的执行中的空间与人员处置”。条目把背景概括为“隗嚣死后，其子隗纯降汉，陇右主要割据力量被解除”发生的政治与区域条件共同作用，并指出该节点改变后续资源、联盟或制度处置的条件。这个节点进入区域社会时，要经由文书、道路、军队、市场、寺院和家庭等不同中介；因此，政治行动的宣布、地方接收和日常生活的改变不能被视作同一时刻完成。

本条依据《后汉书·隗嚣传》；《后汉纪》卷八；《资治通鉴》卷四十三，证据提示为“核心事项已有既有账本来源；本分解节点的参与者、执行范围和时间先后仍须逐案核验”。这些材料可支持年序、官方决定、战事或特定地点的线索，却难以直接统计所有人的财富、迁徙、信仰和动机。更稳妥的读法是区分诏令与实施、军报与人口、行政称谓与自我认同，并让地方材料的缺环限制解释的范围。

由此可见，该事件不是孤立的年表点，而是资源、信息与权力关系重新配置的一环。不同地区的官员、社群和家庭可能以不同方式承担征发、保护、迁移或宗教活动；现存来源只能照亮其中一部分，不能替沉默者制造统一的经验。

\eventanchor{EH-LEDGER-006-5}{建武十年（34年）·“隗嚣死后，其子隗纯降汉，陇右主要割据力量被解除”的后续制度或军政调整}账本记录东汉与陇右集团在建武十年（34年）发生““隗嚣死后，其子隗纯降汉，陇右主要割据力量被解除”的后续制度或军政调整”。条目把背景概括为核心节点发生后仍须处理其遗留问题，并指出后续影响须在不同地区和材料类型中分别核验。这个节点进入区域社会时，要经由文书、道路、军队、市场、寺院和家庭等不同中介；因此，政治行动的宣布、地方接收和日常生活的改变不能被视作同一时刻完成。

本条依据《后汉书·隗嚣传》；《后汉纪》卷八；《资治通鉴》卷四十三，证据提示为“核心事项已有既有账本来源；本分解节点的参与者、执行范围和时间先后仍须逐案核验”。这些材料可支持年序、官方决定、战事或特定地点的线索，却难以直接统计所有人的财富、迁徙、信仰和动机。更稳妥的读法是区分诏令与实施、军报与人口、行政称谓与自我认同，并让地方材料的缺环限制解释的范围。

由此可见，该事件不是孤立的年表点，而是资源、信息与权力关系重新配置的一环。不同地区的官员、社群和家庭可能以不同方式承担征发、保护、迁移或宗教活动；现存来源只能照亮其中一部分，不能替沉默者制造统一的经验。

\eventanchor{EH-LEDGER-006-6}{建武十年（34年）·“隗嚣死后，其子隗纯降汉，陇右主要割据力量被解除”的异源材料与影响范围的复核}账本记录东汉与陇右集团在建武十年（34年）发生““隗嚣死后，其子隗纯降汉，陇右主要割据力量被解除”的异源材料与影响范围的复核”。条目把背景概括为核心节点发生后仍须处理其遗留问题，并指出后续影响须在不同地区和材料类型中分别核验。这个节点进入区域社会时，要经由文书、道路、军队、市场、寺院和家庭等不同中介；因此，政治行动的宣布、地方接收和日常生活的改变不能被视作同一时刻完成。

本条依据《后汉书·隗嚣传》；《后汉纪》卷八；《资治通鉴》卷四十三，证据提示为“核心事项已有既有账本来源；本分解节点的参与者、执行范围和时间先后仍须逐案核验”。这些材料可支持年序、官方决定、战事或特定地点的线索，却难以直接统计所有人的财富、迁徙、信仰和动机。更稳妥的读法是区分诏令与实施、军报与人口、行政称谓与自我认同，并让地方材料的缺环限制解释的范围。

由此可见，该事件不是孤立的年表点，而是资源、信息与权力关系重新配置的一环。不同地区的官员、社群和家庭可能以不同方式承担征发、保护、迁移或宗教活动；现存来源只能照亮其中一部分，不能替沉默者制造统一的经验。

\eventanchor{EH-LEDGER-007-1}{建武十一年（35年）·“光武帝下诏限制将战乱中掠卖者继续为奴婢，并令核实释放”的前期动员与资源准备}账本记录东汉在建武十一年（35年）发生““光武帝下诏限制将战乱中掠卖者继续为奴婢，并令核实释放”的前期动员与资源准备”。条目把背景概括为当时的权力、资源与信息约束推动相关过程展开，并指出这一过程为“光武帝下诏限制将战乱中掠卖者继续为奴婢，并令核实释放”提供可观察的条件。这个节点进入区域社会时，要经由文书、道路、军队、市场、寺院和家庭等不同中介；因此，政治行动的宣布、地方接收和日常生活的改变不能被视作同一时刻完成。

本条依据《后汉书·光武帝纪下》；《后汉书·刘隆传》；《后汉纪》卷九，证据提示为“核心事项已有既有账本来源；本分解节点的参与者、执行范围和时间先后仍须逐案核验”。这些材料可支持年序、官方决定、战事或特定地点的线索，却难以直接统计所有人的财富、迁徙、信仰和动机。更稳妥的读法是区分诏令与实施、军报与人口、行政称谓与自我认同，并让地方材料的缺环限制解释的范围。

由此可见，该事件不是孤立的年表点，而是资源、信息与权力关系重新配置的一环。不同地区的官员、社群和家庭可能以不同方式承担征发、保护、迁移或宗教活动；现存来源只能照亮其中一部分，不能替沉默者制造统一的经验。

\eventanchor{EH-LEDGER-007-2}{建武十一年（35年）·“光武帝下诏限制将战乱中掠卖者继续为奴婢，并令核实释放”的命令、联盟或地方响应}账本记录东汉在建武十一年（35年）发生““光武帝下诏限制将战乱中掠卖者继续为奴婢，并令核实释放”的命令、联盟或地方响应”。条目把背景概括为当时的权力、资源与信息约束推动相关过程展开，并指出这一过程为“光武帝下诏限制将战乱中掠卖者继续为奴婢，并令核实释放”提供可观察的条件。这个节点进入区域社会时，要经由文书、道路、军队、市场、寺院和家庭等不同中介；因此，政治行动的宣布、地方接收和日常生活的改变不能被视作同一时刻完成。

本条依据《后汉书·光武帝纪下》；《后汉书·刘隆传》；《后汉纪》卷九，证据提示为“核心事项已有既有账本来源；本分解节点的参与者、执行范围和时间先后仍须逐案核验”。这些材料可支持年序、官方决定、战事或特定地点的线索，却难以直接统计所有人的财富、迁徙、信仰和动机。更稳妥的读法是区分诏令与实施、军报与人口、行政称谓与自我认同，并让地方材料的缺环限制解释的范围。

由此可见，该事件不是孤立的年表点，而是资源、信息与权力关系重新配置的一环。不同地区的官员、社群和家庭可能以不同方式承担征发、保护、迁移或宗教活动；现存来源只能照亮其中一部分，不能替沉默者制造统一的经验。

\eventanchor{EH-LEDGER-007-3}{建武十一年（35年）·光武帝下诏限制将战乱中掠卖者继续为奴婢，并令核实释放}账本记录东汉在建武十一年（35年）发生“光武帝下诏限制将战乱中掠卖者继续为奴婢，并令核实释放”。条目把背景概括为“光武帝下诏限制将战乱中掠卖者继续为奴婢，并令核实释放”发生的政治与区域条件共同作用，并指出该节点改变后续资源、联盟或制度处置的条件。这个节点进入区域社会时，要经由文书、道路、军队、市场、寺院和家庭等不同中介；因此，政治行动的宣布、地方接收和日常生活的改变不能被视作同一时刻完成。

本条依据《后汉书·光武帝纪下》；《后汉书·刘隆传》；《后汉纪》卷九，证据提示为“核心事项已有既有账本来源；本分解节点的参与者、执行范围和时间先后仍须逐案核验”。这些材料可支持年序、官方决定、战事或特定地点的线索，却难以直接统计所有人的财富、迁徙、信仰和动机。更稳妥的读法是区分诏令与实施、军报与人口、行政称谓与自我认同，并让地方材料的缺环限制解释的范围。

由此可见，该事件不是孤立的年表点，而是资源、信息与权力关系重新配置的一环。不同地区的官员、社群和家庭可能以不同方式承担征发、保护、迁移或宗教活动；现存来源只能照亮其中一部分，不能替沉默者制造统一的经验。

\eventanchor{EH-LEDGER-007-4}{建武十一年（35年）·“光武帝下诏限制将战乱中掠卖者继续为奴婢，并令核实释放”的执行中的空间与人员处置}账本记录东汉在建武十一年（35年）发生““光武帝下诏限制将战乱中掠卖者继续为奴婢，并令核实释放”的执行中的空间与人员处置”。条目把背景概括为“光武帝下诏限制将战乱中掠卖者继续为奴婢，并令核实释放”发生的政治与区域条件共同作用，并指出该节点改变后续资源、联盟或制度处置的条件。这个节点进入区域社会时，要经由文书、道路、军队、市场、寺院和家庭等不同中介；因此，政治行动的宣布、地方接收和日常生活的改变不能被视作同一时刻完成。

本条依据《后汉书·光武帝纪下》；《后汉书·刘隆传》；《后汉纪》卷九，证据提示为“核心事项已有既有账本来源；本分解节点的参与者、执行范围和时间先后仍须逐案核验”。这些材料可支持年序、官方决定、战事或特定地点的线索，却难以直接统计所有人的财富、迁徙、信仰和动机。更稳妥的读法是区分诏令与实施、军报与人口、行政称谓与自我认同，并让地方材料的缺环限制解释的范围。

由此可见，该事件不是孤立的年表点，而是资源、信息与权力关系重新配置的一环。不同地区的官员、社群和家庭可能以不同方式承担征发、保护、迁移或宗教活动；现存来源只能照亮其中一部分，不能替沉默者制造统一的经验。

\eventanchor{EH-LEDGER-007-5}{建武十一年（35年）·“光武帝下诏限制将战乱中掠卖者继续为奴婢，并令核实释放”的后续制度或军政调整}账本记录东汉在建武十一年（35年）发生““光武帝下诏限制将战乱中掠卖者继续为奴婢，并令核实释放”的后续制度或军政调整”。条目把背景概括为核心节点发生后仍须处理其遗留问题，并指出后续影响须在不同地区和材料类型中分别核验。这个节点进入区域社会时，要经由文书、道路、军队、市场、寺院和家庭等不同中介；因此，政治行动的宣布、地方接收和日常生活的改变不能被视作同一时刻完成。

本条依据《后汉书·光武帝纪下》；《后汉书·刘隆传》；《后汉纪》卷九，证据提示为“核心事项已有既有账本来源；本分解节点的参与者、执行范围和时间先后仍须逐案核验”。这些材料可支持年序、官方决定、战事或特定地点的线索，却难以直接统计所有人的财富、迁徙、信仰和动机。更稳妥的读法是区分诏令与实施、军报与人口、行政称谓与自我认同，并让地方材料的缺环限制解释的范围。

由此可见，该事件不是孤立的年表点，而是资源、信息与权力关系重新配置的一环。不同地区的官员、社群和家庭可能以不同方式承担征发、保护、迁移或宗教活动；现存来源只能照亮其中一部分，不能替沉默者制造统一的经验。
